\documentclass[11pt]{article}



\usepackage[T1]{fontenc}
\usepackage[utf8]{inputenc}
\usepackage[english]{babel}
\usepackage[margin=1in]{geometry}
\usepackage{lmodern}
\usepackage{microtype}
\usepackage{mathtools}
\usepackage{amsmath,amssymb,amsthm}
\usepackage{array}
\usepackage{booktabs}
\usepackage{tabularx}
\usepackage{longtable}
\usepackage{float}
\usepackage{graphicx}
\usepackage{tikz}
\usetikzlibrary{arrows.meta,calc,fit,positioning}
\usepackage{pdflscape}
\usepackage{caption}
\usepackage{enumitem}
\usepackage{fancyhdr}
\usepackage[numbers]{natbib}
\usepackage{doi}
\usepackage{xurl}
\usepackage{hyperref}
\usepackage{cleveref}
\usepackage{orcidlink}
\usepackage{titling}

\newtheorem{theorem}{Theorem}[section]
\newtheorem{lemma}[theorem]{Lemma}
\newtheorem{proposition}[theorem]{Proposition}
\newtheorem{corollary}[theorem]{Corollary}
\theoremstyle{definition}
\newtheorem{definition}[theorem]{Definition}
\newtheorem{remark}[theorem]{Remark}
\newtheorem{warning}[theorem]{Warning}
\newtheorem{obligation}[theorem]{Obligation}
\newtheorem{nonclaim}[theorem]{Nonclaim}

\urlstyle{same}
\captionsetup{font=small, labelfont=bf, labelsep=period}
\captionsetup[table]{skip=6pt, position=top}
\captionsetup[figure]{skip=6pt}
\crefname{theorem}{theorem}{theorems}
\Crefname{theorem}{Theorem}{Theorems}
\crefname{lemma}{lemma}{lemmas}
\Crefname{lemma}{Lemma}{Lemmas}
\crefname{proposition}{proposition}{propositions}
\Crefname{proposition}{Proposition}{Propositions}
\crefname{corollary}{corollary}{corollaries}
\Crefname{corollary}{Corollary}{Corollaries}
\crefname{definition}{definition}{definitions}
\Crefname{definition}{Definition}{Definitions}
\crefname{remark}{remark}{remarks}
\Crefname{remark}{Remark}{Remarks}
\crefname{warning}{warning}{warnings}
\Crefname{warning}{Warning}{Warnings}
\crefname{obligation}{obligation}{obligations}
\Crefname{obligation}{Obligation}{Obligations}
\crefname{nonclaim}{nonclaim}{nonclaims}
\Crefname{nonclaim}{Nonclaim}{Nonclaims}
\setlist[enumerate]{leftmargin=*, itemsep=2pt, topsep=4pt}
\setlength{\droptitle}{-3.2em}
\pretitle{\begin{center}\LARGE\bfseries}
\posttitle{\par\end{center}\vspace{0.5em}}
\preauthor{\begin{center}\normalsize}
\postauthor{\par\end{center}\vspace{-0.4em}}
\predate{\begin{center}\small}
\postdate{\par\end{center}\vspace{0.6em}}
\clubpenalty=10000
\widowpenalty=10000
\displaywidowpenalty=10000
\interfootnotelinepenalty=10000
\allowdisplaybreaks[2]
\fancypagestyle{firstpage}{\fancyhf{}
  \fancyfoot[C]{\scriptsize Automorph Inc., Wilmington, DE, USA \quad \texttt{ioannis@automorph.io} \quad \textcopyright\ 2026 \quad CC-BY 4.0 \quad \doi{10.5281/zenodo.20713968}}
  \renewcommand{\headrulewidth}{0pt}
  \renewcommand{\footrulewidth}{0pt}
}



\newcommand{\Q}{\mathbb{Q}}
\newcommand{\R}{\mathbb{R}}
\newcommand{\Z}{\mathbb{Z}}
\newcommand{\Qp}{\mathbb{Q}_p}
\newcommand{\Zp}{\mathbb{Z}_p}
\newcommand{\Sha}{\operatorname{Sha}}
\newcommand{\Reg}{\operatorname{Reg}}
\newcommand{\Tam}{\operatorname{Tam}}
\newcommand{\tors}{\operatorname{tors}}
\newcommand{\tr}{\operatorname{tr}}
\newcommand{\Pf}{\operatorname{Pf}}
\newcommand{\Fitt}{\operatorname{Fitt}}
\newcommand{\cores}{\operatorname{cores}}
\newcommand{\rank}{\operatorname{rank}}
\newcommand{\Hf}{H^1_f}
\newcommand{\kapr}{\kappa_r}

\newcommand{\Sel}{\operatorname{Sel}}
\newcommand{\SelBSD}{\mathrm{Sel}^{!}_{\mathrm{BSD}}}   \newcommand{\rhoE}{\rho_E}                              \newcommand{\PiBSD}{\Pi_{\mathrm{BSD}}}                 \newcommand{\GammaPD}{\Gamma_{\mathrm{BSD}}^{\mathrm{padic\text{-}descent}}}
\newcommand{\GammaSP}{\Gamma_{\mathrm{BSD}}^{\mathrm{Sha\text{-}persistence}}}
\newcommand{\GammaGZ}{\Gamma_{\mathrm{BSD}}^{\mathrm{higher\text{-}GZ\text{-}fixity}}}
\newcommand{\chiCTp}{\chi_{CT,p}}                       \newcommand{\AofE}{\mathcal{A}_E}                        \newcommand{\etap}{\eta_p}                              \newcommand{\RefStable}{\mathbf{RefStable\,AOR}}

\setlength{\emergencystretch}{3.5em}
\sloppy

\title{A Conditional Proof of the Strong Birch--Swinnerton-Dyer Conjecture}
\author{Ioannis Tsiokos\,\orcidlink{0009-0009-7659-5964}}
\date{31 August 2026}
\hypersetup{
  pdftitle={A Conditional Proof of the Strong Birch--Swinnerton-Dyer Conjecture},
  pdfauthor={Ioannis Tsiokos},
  hidelinks
}

\begin{document}
\maketitle
\thispagestyle{firstpage}

\begin{abstract}
We make precise what a conditional closure of Strong
Birch--Swinnerton-Dyer rests on.  The construction is a saturated BSD
trace shell \(\SelBSD\) whose scalar residual \(\rhoE\), the difference
between the analytic leading coefficient and the Strong-BSD arithmetic
product, vanishes exactly when the Strong-BSD scalar identity holds,
together with a structured readout predicate \(\PiBSD\) whose
readout-level equivalence with that identity is a theorem rather than a
definition.
The deep arithmetic is isolated into three named recognition sources:
a \(p\)-adic descent/Tamagawa correction, a signed-Selmer
\(\Sha\)-persistence source, and a higher-rank Gross--Zagier fixity
source.  These are supplied as typed carriers, not derived.  They are
combined with headline imported theorem stacks: a Cassels--Tate Pfaffian
comparison, a Selmer-complex pairing substrate, and a Beilinson
rank-\(\geq 2\) comparison.  Under the standing shell-closure
assumption, these sources and headline imports make the framework master
theorem force \(\rhoE=0\), and the readout equivalence yields Strong BSD.
Outside the framework this is the conditional theorem that the standing
shell-closure record, the three recognition sources, and the headline
imports imply Strong BSD; it is not an unconditional proof.  The closed
carrier is then read as a
non-eliminative Audited Operational Realisability instance: the
recognition records are bridged rather than discharged, while the
rank-\(\geq 2\) generalized Gross--Zagier identity and additive
\(p\leq 3\) Iwasawa main conjecture remain open.  The architecture is
formalized in Lean as a checked conditional schema, not as an
independent derivation.
\end{abstract}

\medskip
\noindent\textbf{Keywords.} Birch--Swinnerton-Dyer conjecture; Strong BSD;
conditional closure; Cassels--Tate pairing; Selmer complex; Iwasawa main
conjecture; formal verification (Lean); Six Birds Theory; emergence calculus.

\medskip
\noindent
Conventions: $E/\Q$, $M_E=h^1(E)(1)$, $r=\operatorname{rank}E(\Q)$,
$r_{\mathrm{an}}=\operatorname{ord}_{s=1}L(E,s)$. This paper mechanizes the
framework-structural closure architecture only; the deep arithmetic enters as
imported proof-carrying structures and typed recognition-source carriers, never
as Lean axioms. The landing is a conditional closure, not an unconditional proof
of BSD. The \(\kapr\) normalization is reused from
\citet{TsiokosBSDApparatus2026}.

\section{Introduction}\label{sec:introduction}


What would a conditional closure of Strong Birch--Swinnerton-Dyer~\citep{BirchSwinnertonDyer1965,Tate1966bsd} look
like, and what exactly would it rest on?  The classical landscape already
contains powerful conditional and partial routes: Gross--Zagier and
Kolyvagin~\citep{GrossZagier1986,Kolyvagin1989} control the rank-one Heegner setting; Skinner--Urban and
Iwasawa-main-conjecture methods~\citep{SkinnerUrban2014} control large ordinary regions;
Cassels--Tate and Selmer-complex formalisms~\citep{Flach1990,Nekovar2006} control part of the
quadratic and determinant bookkeeping.  These results do not, by
themselves, give a single uniform proof of the strong formula in every
rank and at every bad prime.  The question addressed here is narrower
and more structural: if the remaining arithmetic inputs are isolated
with precise names, can one assemble them into a single audited closure
whose scalar readout is exactly the Strong-BSD identity?

The answer is deliberately conditional.  This is \emph{not} an
unconditional classical proof of BSD, scalar or strong.  The result is
theorem-grade within the Six Birds closure discipline, conditional on
the standing shell-closure record, the named recognition sources, and the
headline imported theorem stacks.  Outside the framework it reads as the
conditional theorem
\[
  \begin{aligned}
  &\SelBSD\text{-closure}
    \wedge \GammaPD \wedge \GammaSP \wedge \GammaGZ
    \wedge \chiCTp\text{-imports}
    \wedge \AofE\text{-imports} \\
  &\qquad {}\wedge\ \text{Beilinson-imports}
    \Longrightarrow \text{Strong BSD}.
  \end{aligned}
\]
This implication is the paper's result in its plainest form.  The work
below does not erase the hypotheses; it makes their roles explicit and
shows that, once they are supplied, the Strong-BSD scalar identity is
forced by a single closure architecture.

The first component of the architecture is a separation of the deep
arithmetic into two kinds of inputs.  The first kind consists of three
recognition sources.  A recognition source is named arithmetic content
carried by the formed BSD layer: it is supplied as part of the closure
object and used as a hypothesis, not derived from the general framework.
The three sources are \(\GammaPD\), the additive-prime \(p\)-adic
descent/Tamagawa correction; \(\GammaSP\), the signed-Selmer
persistence of the \(\Sha\)-factor through local cover change; and
\(\GammaGZ\), the higher-rank Gross--Zagier-type fixity of the
archimedean regulator readout.  They correspond, respectively, to the
local \(p\)-adic factor, the \(\Sha\)-persistence factor, and the
rank-\(\geq 2\) regulator/fixity factor in the Strong-BSD product.  In
the formal development they are typed carriers of hypotheses, not
free-standing assumptions of the framework; in the paper they are treated
as named recognition content and are introduced before they are used.

The second kind of input consists of imported theorem stacks.  The three
headline imports are the ones used directly in the landing theorem.  The
\(\chiCTp\) stack packages the Cassels--Tate Pfaffian comparison and its
sign convention.  The \(\AofE\) stack packages the Selmer-complex pairing
substrate needed to compare the determinant formulation with the
Cassels--Tate side.  The Beilinson stack packages the rank-\(\geq 2\)
archimedean determinant--\(L\)-value comparison.  Later sections also
record the overconvergent \(\eta\) and \(T_{\mathrm{CASCADE}}\) stacks as
supporting conditional imports, but those are not hidden inside the three
headline imports.  These imports name classical arithmetic geometry as
provenance: Sakamoto, Macias Castillo--Sano, Nekov\'a\v{r}, Flach, the
Cassels--Tate literature, Burns--Flach, Beilinson, and the surrounding
Selmer-complex and determinant technology.  They are not reproved here.
They enter as proof-carrying assumptions, and later sections keep their
individual roles separate.

The second component is the saturated BSD trace shell.  It is the
five-part object
\[
  \SelBSD=(E,\rhoE,J^{!}_{\mathrm{BSD}},
  \mathrm{Vis}^{!}_{\mathrm{BSD}},\mathrm{Audit}^{!}_{\mathrm{BSD}}),
\]
where \(E/\Q\) is the elliptic curve, \(J^{!}_{\mathrm{BSD}}\) records
the duality or involutive datum, \(\mathrm{Vis}^{!}_{\mathrm{BSD}}\)
records the visible scalar readout, and
\(\mathrm{Audit}^{!}_{\mathrm{BSD}}\) records the closure audit.  The
central scalar is the residual
\[
  \rhoE
  =
  \frac{L^{(r)}(E,1)}{r!}
  -
  \frac{\#\Sha(E/\Q)\,\Reg_{\mathrm{NT}}(E)\,\Omega_E
    \prod_v c_v(E)}
       {|E(\Q)_{\tors}|^2}.
\]
Thus \(\rhoE=0\) is not a new BSD-like statement: it is exactly the
standard Strong-BSD scalar identity written as a residual equation.
The point of the shell is to keep the analytic side, the arithmetic
product, the duality datum, the visible readout, and the audit data in
one typed object, so that later arguments can say precisely which input
changes which part of the scalar package.

The recognition predicate \(\PiBSD\) is the shell's structured readout.
It is not defined to be Strong BSD.  It is a quantified predicate that
asks for the three recognition-source readouts in their proper ranges:
the \(p\)-adic descent readout for every additive bad prime and
trivializing local lift in scope, the signed-Selmer persistence readout
for every admissible signed-Selmer prime, and the higher-Gross--Zagier
fixity readout in rank at least two, with the classical low-rank lane
for ranks \(0\) and \(1\).  The readout-level equivalence says that
vanishing of the shell residual is equivalent to the usual Strong-BSD
scalar identity.  This equivalence is a genuine algebraic step: it
rewrites the residual, rather than making \(\PiBSD\) a synonym for the
conjecture.

The headline landing then has a short conceptual chain.  The recognition
sources and imports are supplied.  The higher-rank fixity source forces
the shell residual to vanish, after the \(\kapr\)-normalization identifies
the cascade regulator product with the standard
\(\#\Sha/|E(\Q)_{\tors}|^2\) factor.  The readout equivalence converts
\(\rhoE=0\) into
\[
  \frac{L^{(r)}(E,1)}{r!}
  =
  \frac{\#\Sha(E/\Q)\,\Reg_{\mathrm{NT}}(E)\,\Omega_E
    \prod_v c_v(E)}
       {|E(\Q)_{\tors}|^2}.
\]
The master-theorem applicability audit checks that the three recognition
sources have the correct framework profile, and the composite signature
checks that the three source-governed factors are not disposable: changing
the \(\Sha\)-factor, the regulator factor, or the Tamagawa factor changes
the scalar package.  This is the sense in which the conclusion is a
closure result rather than a restatement of its target.

Two ingredients are imported from \citet{TsiokosBSDApparatus2026}.
First, the \(\kapr\) normalization in that work is the bridge from the cascade leading-term
form to the standard Strong-BSD factor \(\#\Sha/|E(\Q)_{\tors}|^2\).
Second, its adequacy and five-column framing provide the
diagnostic language for separating height, period, finite, \(p\)-adic,
and determinant contributions.  This article uses those results as cited
inputs; it does not refer to internal labels of the cited work.  The AOR
terminology used at the end of the paper is the Audited Operational
Realisability language of \cite{TsiokosAOR2026}, and the recognition-mode
discipline is part of the audited interaction calculus of
\cite{TsiokosFoundationsIII2026}.

The final move is a reading of the completed carrier as a refinement-stable
audited object.  Audited Operational Realisability is a refinement-stable,
audit-closed membership criterion: informally, an object qualifies when
the residuals left by refinement have been accounted for by approved
mechanical records, imported records, or bridged recognition records.  In
this paper the object is \(\SelBSD\).  The resulting statement is the
framework's ``Strong BSD holds in reality'' reading, but it is
non-eliminative: the recognition records are reclassified as bridged
records, not made to disappear.

We close the introduction by fixing the scope.  This is not an
unconditional classical proof of BSD.  It is a conditional closure under
the standing shell-closure record, the three recognition sources, and the
headline imported theorem stacks displayed above.  The rank-\(\geq 2\)
generalized Gross--Zagier identity
\(T_{E12}\) and the additive \(p\leq 3\) Iwasawa main conjecture
\(T_{\mathrm{BAD}}\) remain open audit-result no-gos: the paper records
where they are needed and how the recognition sources carry them, but
does not prove them.  Likewise, the AOR-instance reading is
non-eliminative; it records an audit-closed presentation of the conditional
closure, not the discharge of every open arithmetic obligation.
 \section{The imported theorem stacks}
\label{sec:imports}

The closure separates its inputs into two classes.  The present section
records the imported theorem stacks: established arithmetic geometry
that the closure consumes as proof-carrying hypotheses.  The next section
records the recognition sources: arithmetic content carried by the
formed BSD layer.  This distinction matters.  The imported stacks are not
new theorems here, and the closure never derives them
from the framework.  It derives consequences from them, after their
provenance and downstream role have been made explicit.

There are three headline imports in this section.  The first controls the
Cassels--Tate side of the finite \(p\)-primary comparison.  The second
supplies the Selmer-complex pairing substrate \( \AofE \).  The third
supplies the rank-\(\geq 2\) archimedean determinant comparison.  Together
they keep the classical arithmetic geometry used by the landing theorem out
of the framework argument while still making clear exactly which classical
inputs are being used.  The overconvergent \(\eta\) and
\(T_{\mathrm{CASCADE}}\) imports are introduced later as supporting
conditional stacks and are counted separately in App~\ref{app:formalization}.

The symbol \(\chiCTp\) refers here to a Cassels--Tate comparison: a
comparison between a square-root Fitting object on a Selmer-complex
cohomology group and the Pfaffian scalar attached to the Cassels--Tate
pairing.  In ordinary terms, this is the input that lets the \(p\)-primary
finite piece be read with the correct orientation and sign.

\begin{definition}[Imported $\chi_{CT,p}$ theorem stack]\label{def:closure:chi-ct-p-import}
A proof-carrying structure bundling the eight named external inputs
\(T_{E1}\ldots T_{E8}\): Sakamoto's Stark-system theory~\citep{Sakamoto2018};
the determinant/Stark-system and derived-height comparisons of Macias
Castillo--Sano~\citep[Thms.~3.4 and~4.2]{MaciasCastilloSano2026};
Knudsen--Mumford determinant theory~\citep{KnudsenMumford1976};
Nekov\'a\v{r}'s Selmer-complex pairing~\citep[\S10.8.7]{Nekovar2006};
Flach's pairing comparison~\citep{Flach1990}; and the classical
2-Selmer normalization of Fisher--Schaefer--Stoll~\citep[Thm.~6.3]{FisherSchaeferStoll2010}.
The results are stored as proposition fields,
plus the unconditional sign closure
\(\Sigma_{\mathrm{NekCT}}=+1\).  It is taken as a hypothesis by
\Cref{thm:closure:chi-ct-p-comparison}; it is not derived in this
paper.\footnote{Encoded in Lean as a typed structure carrier
\texttt{SixBirdsBSD.\allowbreak Closure.\allowbreak Imports.\allowbreak chiCTpImport}
bundling the imported theorems as hypothesis fields; not a Lean axiom.
See App~\ref{app:formalization}.}
\end{definition}

The downstream role of this stack is narrow.  It is not a general
replacement for the Cassels--Tate literature.  It is the specific
orientation-respecting package needed for the \(\chiCTp\) theorem:
given the usual finite \(\Sha[p^\infty]\), nondegeneracy, and
Selmer-complex hypotheses, it supplies the Pfaffian comparison and the
sign closure \(\Sigma_{\mathrm{NekCT}}=+1\).

The second import packages the Selmer-complex pairing substrate denoted
\(\AofE\).  This notation is local to this article: \(\AofE\) is the
unified Selmer-complex pairing object, not the analytic leading coefficient
written \(A_E\) in \citet{TsiokosBSDApparatus2026}.  Mathematically, it is the
ambient category, Selmer complex, and pairing target in which the
Cassels--Tate pairing is intrinsic and compatible with the
determinant-line formulation.

\begin{definition}[Imported $\AofE$ Selmer-complex pairing]\label{def:closure:a-e-import}
A proof-carrying structure for the \(\AofE\) unified Selmer-complex
substrate (Nekov\'a\v{r} Selmer complex~\citep{Nekovar2006} with intrinsic Cassels--Tate
pairing; Burns--Flach and related Macias Castillo--Sano--Nekov\'a\v{r}
ETNC-adjacent input~\citep{BurnsFlach2001,BlochKato1990,FontainePerrinRiou1994,BurnsMaciasCastilloSeo2024,MaciasCastilloSano2026}
\(\AofE^{\mathrm{Sel}}\)): a typed target category with a
bilinear/categorical pairing field.  It is used only as a
hypothesis.\footnote{Encoded in Lean as a typed structure carrier
\texttt{SixBirdsBSD.\allowbreak Closure.\allowbreak Imports.\allowbreak aEImport}
bundling the imported theorem as a hypothesis field; not a Lean axiom.
See App~\ref{app:formalization}.}
\end{definition}

This import is the bridge between two ways of seeing the same finite
arithmetic information.  On one side is the Selmer-complex description
with its intrinsic Cassels--Tate pairing.  On the other is the
determinant-line language used in the closure landing.  The \(\AofE\)
carrier records that the target category and pairing data have already
been supplied, so later arguments may use the pairing without
reconstructing the Selmer-complex theory.

The third import is archimedean.  In rank at least two, the landing
needs a comparison between the Beilinson determinant side and the
leading \(L\)-value side.  This is where the regulator-to-determinant
identification enters.  The closure treats it as a classical imported
comparison, separate from the higher-rank recognition source that will
encode the remaining Gross--Zagier-type fixity.

\begin{definition}[Imported Beilinson rank-$\geq 2$ comparison]\label{def:closure:beilinson-import}
A proof-carrying structure for the Beilinson determinant/\(L\)-value~\citep{Beilinson1985}
comparison in rank \(r\geq 2\): the archimedean
regulator-to-determinant-line identification used by the rank-\(\geq 2\)
landing.  It is used only as a hypothesis.\footnote{Encoded in Lean as a
typed structure carrier
\texttt{SixBirdsBSD.\allowbreak Closure.\allowbreak Imports.\allowbreak beilinsonImport}
bundling the imported theorem as a hypothesis field; not a Lean axiom.
See App~\ref{app:formalization}.}
\end{definition}

The point of isolating this comparison is to prevent the closure
argument from hiding an archimedean theorem inside the recognition
predicate.  The Beilinson import supplies the determinant/\(L\)-value
identification; the recognition source introduced later supplies the
formed-layer fixity condition that makes the rank-\(\geq 2\) readout
land on the Strong-BSD scalar.

\begin{nonclaim}\label{nonclaim:closure:imports-not-derived}
The \(\chiCTp\) comparison remains conditional on its imported stack
\(T_{E1}\ldots T_{E8}\); the import structures are sourced assumptions,
not results of this paper.  The closure derives consequences
\emph{from} them and never proves them.  They are carried as typed
hypothesis structures and introduce no project-local axioms.
\end{nonclaim}

 \section{The three recognition sources}
\label{sec:recognition-sources}


This section names the three recognition sources used by the closure.
A recognition source is named arithmetic content on which the closure
rests: it is supplied to the construction, not derived from the general
framework, and it is carried as a typed hypothesis.  The concrete
example to keep in mind is the \(p\)-adic descent congruence below,
which says how a local descent comparison repairs the Tamagawa
coarsening at an additive prime.

The three sources govern the three factors that are not supplied by the
formal shell alone.  The source \(\GammaPD\) governs the local
\(p\)-adic/Tamagawa correction in the odd-additive scope of its definition.
The source \(\GammaSP\) governs the
persistence of the \(\Sha\)-factor through signed-Selmer cover change.
The source \(\GammaGZ\) governs the rank-\(\geq 2\)
regulator/fixity readout.  Together they are the arithmetic content that
the conditional landing will consume.
\Cref{tab:recognition-sources} records the same three sources as a compact
audit table.
\begin{table}[tbp]
\centering
\footnotesize
\caption{The three recognition sources consumed by the conditional landing.}
\label{tab:recognition-sources}
\begin{tabularx}{\textwidth}{@{}>{\raggedright\arraybackslash}p{0.18\textwidth}
>{\raggedright\arraybackslash}p{0.16\textwidth}
>{\raggedright\arraybackslash}p{0.34\textwidth}
>{\raggedright\arraybackslash}p{0.15\textwidth}
>{\raggedright\arraybackslash}X@{}}
\toprule
Source & Mode-B residual & Arithmetic content & Foundations-IV & Stage-6 \\
\midrule
\(\GammaPD\) & \(D_3'''\) & Odd-additive-prime \(p\)-adic descent congruence
with the Tamagawa correction \(c_p(E)^{-1}\) and a route unit
\(\varepsilon_p(E,L)\). & F2/F3/F4 & FP \(\to\) (ii) \\
\addlinespace
\(\GammaSP\) & \(D_2'''\) & Signed-Selmer cover-change determinant persists
as the local component factor \(c_p(E)\) modulo \(\Zp^\times\). &
F19/F21/F29 & SAU \(\to\) (iii) \\
\addlinespace
\(\GammaGZ\) & \(D_2''\) & Higher-rank fixity
\(\psi_-(E,r)=0\), identifying the analytic leading coefficient with the
\(\kapr\)-normalized regulator product. & F14/F27/F40 & VDE \(\to\) (ii) \\
\bottomrule
\end{tabularx}
\end{table}

\begin{definition}[$\GammaPD$ p-adic descent carrier]\label{def:closure:gamma-padic-descent}
For \(E/\Q\), an odd additive prime \(p\)~\citep{Silverman1994} with shallow Kodaira type
II/III, \(c_p(E)\in\{1,2\}\), \(p\nmid c_p(E)\),
\(\Sha(E/\Q)[p^\infty]=0\), and a finite tame
inertia-trivializing extension \(L/\Qp\), the formed local layer carries
an ordinary/anticyclotomic descent comparison
\[
  \mathrm{descent}_p^{\mathrm{OC}}:
  LC_{\mathrm{triv}}(L_p^{\mathrm{OC}}(E_L))
  \longrightarrow
  LC_{\mathrm{triv}}(L_p^{\mathrm{OC}}(E))
\]
and a unit \(\varepsilon_p(E,L)\in\Zp^\times\) with
\[
  LC_{\mathrm{triv}}(L_p^{\mathrm{OC}}(E)) \equiv
  c_p(E)^{-1}\,\varepsilon_p(E,L)\,
  \mathrm{descent}_p^{\mathrm{OC}}
    (LC_{\mathrm{triv}}(L_p^{\mathrm{OC}}(E_L)))
  \pmod{\Zp^\times}.
\]
This congruence recognizes that the descent comparison repairs the
Tamagawa coarsening through the factor \(c_p(E)^{-1}\), while the unit
\(\varepsilon_p(E,L)\) records the residue introduced by changing the
local cover.  The arithmetic assertion is used as a hypothesis in the
closure.\footnote{Supplied as a recognition source, encoded in Lean as a
typed structure carrier
\texttt{SixBirdsBSD.\allowbreak Closure.\allowbreak RecognitionSources.\allowbreak gammaPadicDescent}
over a narrowed structural representation; its arithmetic content is a
hypothesis field, not a Lean axiom. See App~\ref{app:formalization}.}
\end{definition}

The second source records persistence of the component measured by
\(\Sha\).  A signed local pairing~\citep{Kobayashi2003,Pollack2003} changes when one passes to a tame
inertia-trivializing extension and then corestricts back.  The source
\(\GammaSP\) says that the determinant of this cover-change operation
is exactly the local component factor modulo a \(p\)-adic unit.

\begin{definition}[$\GammaSP$ signed-Selmer carrier]\label{def:closure:gamma-sha-persistence}
For \(E/\Q\), a signed-Selmer-admissible additive prime \(p\), and a tame
inertia-trivializing extension \(L/\Qp\), the signed-Selmer persistence
layer carries the cover-change determinant, where
\(\mathrm{sp}_p(E,L)\) denotes the signed local pairing,
\[
  C_{\mathrm{loc}}^{\mathrm{III}}(E,p)
  :=
  \det\!\bigl(\cores_{L/\Qp}\circ \mathrm{sp}_p(E,L)\bigr)
\]
and asserts
\[
  C_{\mathrm{loc}}^{\mathrm{III}}(E,p)
  \equiv c_p(E)\pmod{\Zp^\times}.
\]
Thus the signed-Selmer local determinant persists as the Tamagawa
component \(c_p(E)\), up to the unit ambiguity natural in the
\(p\)-adic comparison.\footnote{Supplied as a recognition source,
encoded in Lean as a typed structure carrier
\texttt{SixBirdsBSD.\allowbreak Closure.\allowbreak RecognitionSources.\allowbreak gammaShaPersistence}
over a narrowed structural representation; its arithmetic content is a
hypothesis field, not a Lean axiom. See App~\ref{app:formalization}.}
\end{definition}

The third source is the higher-rank fixity source.  It measures the
difference between the analytic leading coefficient and the arithmetic
regulator product.  The anti-invariant readout is simply
``analytic minus arithmetic''.  In the arithmetic factor, the regulator
is the N\'eron--Tate Gram determinant
\(\det\langle P_i,P_j\rangle\), understood over the
\(\mathrm{GL}_r(\Z)\)-orbit of choices of Mordell--Weil basis.

\begin{definition}[$\GammaGZ$ higher-rank fixity carrier]\label{def:closure:gamma-higher-gz-fixity}
For non-CM \(E/\Q\) of analytic rank \(r\geq 2\), the formed
rank-\(r\) arithmetic/analytic duality layer has anti-invariant readout
\[
  \psi_-(E,r)
  =
  \frac{L^{(r)}(E,1)}{r!}
  -
  \kapr(E)\,\Reg_{\mathrm{NT}}(E)\,\Omega_E\,\Tam(E),
\]
and asserts the higher-Gross--Zagier fixity
\[
  \psi_-(E,r)=0.
\]
Equivalently, the analytic leading coefficient equals the rank-\(r\)
arithmetic product built from \(\kapr(E)\), the N\'eron--Tate regulator,
the period, and the Tamagawa product.  This assertion is used as a
recognition source in the landing.\footnote{Supplied as a recognition source, encoded in Lean as a
typed structure carrier
\texttt{SixBirdsBSD.\allowbreak Closure.\allowbreak RecognitionSources.\allowbreak gammaHigherGZFixity}
over a narrowed structural representation; its arithmetic content is a
hypothesis field, not a Lean axiom. See App~\ref{app:formalization}.}
\end{definition}

This is the source that directly meets the shell residual in high rank.
By the \(\kapr\) normalization of \citet{TsiokosBSDApparatus2026},
the normalization factor identifies
\[
  \kapr(E)=\frac{\#\Sha(E/\Q)}{|E(\Q)_{\tors}|^2}.
\]
Consequently
\[
  \kapr(E)\,\Reg_{\mathrm{NT}}(E)\,\Omega_E\,\Tam(E)
  =
  \frac{\#\Sha(E/\Q)\,\Reg_{\mathrm{NT}}(E)\,\Omega_E
    \prod_v c_v(E)}
       {|E(\Q)_{\tors}|^2},
\]
after reading \(\Tam(E)\) as the Tamagawa product.  Thus
\(\GammaGZ\), through the assertion \(\psi_-(E,r)=0\), is precisely the
higher-rank fixity that will force the residual \(\rhoE\) to vanish in
the landing theorem.

\begin{nonclaim}\label{nonclaim:closure:recognition-sources-not-derived}
No recognition source is derived from framework primitives.  Each
\(\Gamma_{\mathrm{BSD}}^{*}\) is closure content of the formed BSD
layer, encoded as a typed structure carrier whose arithmetic content is
a hypothesis field, not a Lean axiom.
\end{nonclaim}

 \section{The saturated trace shell and its readout}
\label{sec:shell-readout}

The preceding section isolated the three arithmetic recognition sources.
We now assemble the object on which they act.  The saturated BSD trace
shell is a typed carrier for the Strong-BSD scalar comparison: it keeps
the analytic leading coefficient, the arithmetic product, and the audit
data in one place, so that the closure can record which input governs
which scalar factor.  Here ``trace shell'' means a formed mathematical
object together with its visible scalar readout and the bookkeeping
needed to prevent the readout from being a tautological restatement of
the desired identity.

\begin{definition}[Saturated BSD trace shell]\label{def:closure:sel-bsd-shell}
\(\SelBSD\) is the five-tuple
\[
  \SelBSD
  =
  \bigl(E,\ \rhoE,\ J^{!}_{\mathrm{BSD}},\
    \mathrm{Vis}^{!}_{\mathrm{BSD}},\
    \mathrm{Audit}^{!}_{\mathrm{BSD}}\bigr).
\]
Here \(E/\Q\) is the elliptic curve and \(\rhoE\) is the scalar
residual
\[
  \rhoE
  =
  \frac{L^{(r)}(E,1)}{r!}
  -
  \frac{\#\Sha(E/\Q)\,\Reg_{\mathrm{NT}}(E)\,\Omega_E\,\Tam(E)}
       {|E(\Q)_{\tors}|^2}.
\]
Thus \(\rhoE=0\) is exactly the Strong-BSD scalar identity, with
\(\Tam(E)\) read as the Tamagawa product.  The datum
\(J^{!}_{\mathrm{BSD}}\) records the involutive or duality structure
that compares the analytic and arithmetic sides.  The datum
\(\mathrm{Vis}^{!}_{\mathrm{BSD}}\) records the visible scalar readout
of the shell.  The datum \(\mathrm{Audit}^{!}_{\mathrm{BSD}}\) records
the admissibility, anti-tautology, and lawfulness information attached
to the formed object.  The standing closure assumption used here is the
standard assumption that the formed object carries this fixed content
itself: admissibility, anti-tautology, and lawfulness data are part of
the shell rather than consequences inferred after the fact.
\end{definition}

\Cref{fig:sel-bsd-shell-schematic,tab:sel-bsd-shell-fields} summarize the
shell fields and the residual readout used below.
\begin{figure}[tbp]
\centering
\begin{tikzpicture}[
  font=\small,
  box/.style={draw, rounded corners=2pt, align=center, minimum width=2.5cm,
    minimum height=0.85cm, fill=black!3},
  shell/.style={draw, rounded corners=3pt, align=center, minimum width=4.2cm,
    minimum height=1.1cm, fill=black!8},
  arrow/.style={-{Latex[length=2mm]}, thick}
]
\node[shell] (shell) at (0,0) {\(\SelBSD=(E,\rhoE,J^{!},\mathrm{Vis}^{!},\mathrm{Audit}^{!})\)};
\node[box] (analytic) at (-4.6,1.7) {analytic\\leading term};
\node[box] (arith) at (0,1.7) {arithmetic\\product};
\node[box] (audit) at (4.6,1.7) {audit\\records};
\node[box] (readout) at (0,-1.8) {\(\rhoE=0\)\\readout};

\draw[arrow] (analytic) -- (shell);
\draw[arrow] (arith) -- (shell);
\draw[arrow] (audit) -- (shell);
\draw[arrow] (shell) -- (readout);
\node[align=center, font=\footnotesize] at (0,-2.75)
  {the readout is equivalent to the Strong-BSD scalar identity};
\end{tikzpicture}
\caption{The saturated BSD trace shell packages the residual, readout data, and
audit records used by the conditional closure.}
\label{fig:sel-bsd-shell-schematic}
\end{figure}
 \begin{table}[tbp]
\centering
\small
\caption{Fields of the saturated BSD trace shell.}
\label{tab:sel-bsd-shell-fields}
\begin{tabularx}{\textwidth}{@{}>{\raggedright\arraybackslash}p{0.18\textwidth}
>{\raggedright\arraybackslash}p{0.34\textwidth}
>{\raggedright\arraybackslash}X@{}}
\toprule
Field & Role & Paper-prose name \\
\midrule
\(E\) & The elliptic curve over \(\Q\). & base curve \\
\(\rhoE\) & Difference between the analytic leading coefficient and the
Strong-BSD arithmetic product. & scalar residual \\
\(J^{!}_{\mathrm{BSD}}\) & Involutive or duality datum comparing analytic and
arithmetic sides. & duality datum \\
\(\mathrm{Vis}^{!}_{\mathrm{BSD}}\) & Visible scalar readout data attached
to the shell. & visible readout \\
\(\mathrm{Audit}^{!}_{\mathrm{BSD}}\) & Admissibility, anti-tautology, and
lawfulness records. & audit record \\
\addlinespace
\(\rhoE=0\) & The residual vanishing equation, equivalent to the Strong-BSD
scalar identity at readout level. & landing readout \\
\bottomrule
\end{tabularx}
\end{table}

The predicate placed on the shell is not the Strong-BSD identity by
definition.  It is a structured readout predicate: it asks for the three
recognition-source statements, quantified over the local and rank data
where they apply.  This distinction is essential.  The next theorem can
then identify the readout with the Strong-BSD scalar identity; it is not
merely unfolding a definition.

\begin{definition}[$\Pi_{\mathrm{BSD}}$ predicate]\label{def:closure:pi-bsd}
\(\PiBSD(E)\) is the conjunction of the three recognition-source
readouts on \(\SelBSD\), with the relevant quantifiers retained.  First,
for every additive bad prime \(p\) in scope and every tame
inertia-trivializing local lift \(L/\Qp\) in scope, the
\(\GammaPD\) readout supplies the \(p\)-adic descent/Tamagawa
correction.  Second, for every signed-Selmer-admissible prime \(p\),
the \(\GammaSP\) readout supplies the signed-Selmer local determinant
matching.  Third, for analytic rank \(r\geq 2\), the \(\GammaGZ\)
readout supplies the higher-rank fixity assertion
\(\psi_-(E,r)=0\); for \(r\in\{0,1\}\), the predicate uses the
classical low-rank lane.  In particular, \(\PiBSD\) is a structured
quantified predicate over recognition readouts, not the definition
\(\PiBSD:=\mathrm{StrongBSD}\).
\end{definition}

\begin{theorem}[Readout-level equivalence]\label{thm:closure:pi-bsd-iff-strong-bsd}
The scalar readout of \(\PiBSD\) on \(\SelBSD\) is equivalent to the
standard Strong-BSD scalar identity
\[
  \frac{L^{(r)}(E,1)}{r!}
  =
  \frac{\#\Sha(E/\Q)\,\Reg_{\mathrm{NT}}(E)\,\Omega_E\,
    \prod_v c_v(E)}
       {|E(\Q)_{\tors}|^2}.
\]
This is a readout-level equivalence: it identifies the visible scalar
readout with the Strong-BSD scalar identity, but it does not define
\(\PiBSD\) to be Strong BSD and it does not manufacture the recognition
sources from the scalar identity alone.\footnote{Verified in Lean as
\texttt{SixBirdsBSD.\allowbreak Closure.\allowbreak SelShell.\allowbreak piBSDIffStrongBSD};
see App~\ref{app:formalization}.}
\end{theorem}

\begin{proof}
By definition of the shell, the visible scalar residual is
\[
  \rhoE
  =
  \frac{L^{(r)}(E,1)}{r!}
  -
  \frac{\#\Sha(E/\Q)\,\Reg_{\mathrm{NT}}(E)\,\Omega_E\,\Tam(E)}
       {|E(\Q)_{\tors}|^2}.
\]
The readout predicate \(\PiBSD\) supplies the three pieces entering this
residual.  The higher-rank source \(\GammaGZ\) gives
\(\psi_-(E,r)=0\).  After applying the \(\kapr\) normalization of
\citet{TsiokosBSDApparatus2026},
\[
  \kapr(E)=\frac{\#\Sha(E/\Q)}{|E(\Q)_{\tors}|^2},
  \qquad
  \kapr(E)\,\Reg_{\mathrm{NT}}(E)\,\Omega_E\,\Tam(E)
  =
  \frac{\#\Sha(E/\Q)\,\Reg_{\mathrm{NT}}(E)\,\Omega_E
    \prod_v c_v(E)}
       {|E(\Q)_{\tors}|^2}.
\]
Thus the equation \(\psi_-(E,r)=0\) is exactly the assertion
\(\rhoE=0\) in rank at least two.  The source \(\GammaPD\) supplies the
local \(p\)-adic/Tamagawa correction entering the Tamagawa product, and
\(\GammaSP\) supplies the signed-Selmer persistence of the \(\Sha\)
factor; these readouts fix the remaining scalar factors in the same
residual.  Therefore the forward scalar readout of \(\PiBSD\) is
\(\rhoE=0\).

Finally, the equation \(\rhoE=0\) is algebraically the displayed
Strong-BSD scalar identity: moving the arithmetic product to the other
side gives the identity, and subtracting the right-hand side from the
left gives \(\rhoE=0\).  Conversely, if the displayed scalar identity is
already known as a scalar readout, then the same subtraction gives
\(\rhoE=0\).  This reverse implication is only an equivalence of
scalar readouts; it does not construct any of the three
\(\Gamma_{\mathrm{BSD}}^{*}\) predicates.
\end{proof}

\begin{theorem}[Master-theorem applicability]\label{thm:closure:master-theorem-applicability}
The closure master theorem applies to \(\SelBSD\): for the audited
interaction calculus \cite{TsiokosFoundationsIII2026}, correctly
instantiated recognition sources that exhaust the residual factors, and
that pass the no-smuggling audit, force the residual readout
\(\rhoE=0\).\footnote{Verified in Lean as
\texttt{SixBirdsBSD.\allowbreak Closure.\allowbreak SelShell.\allowbreak masterTheoremApplicability};
see App~\ref{app:formalization}.}
\end{theorem}

\begin{proof}
We verify the two applicability requirements for the shell.  First, the
recognition sources are instantiated on exactly the scalar factors of
\(\rhoE\).  The source \(\GammaPD\) governs the local
\(p\)-adic/Tamagawa correction: it supplies the descent congruence and
the Tamagawa factor \(c_p(E)\) up to a \(p\)-adic unit.  The source
\(\GammaSP\) governs the persistence of the \(\Sha\)-factor under the
signed-Selmer cover change: its determinant comparison is the local
component entering the same arithmetic product.  The source
\(\GammaGZ\) governs the rank-\(r\) regulator/fixity factor: it states
that the analytic leading coefficient equals the arithmetic regulator
product after the \(\kapr\) normalization.  These three pieces account
for the local \(p\)-adic/Tamagawa factor, the \(\Sha\)-persistence
factor, and the regulator/fixity factor respectively.  No scalar factor
of \(\rhoE\) remains outside this list, and no factor is assigned to two
different sources.

Second, the no-smuggling condition holds.  The imported theorem stacks
enter only as hypotheses, and the recognition sources enter only as
carried readouts.  The closure does not silently turn an open
arithmetic obligation into a derivation target; it records the
obligation as supplied content and then derives consequences from it.
Thus the recognition-source instantiation is exhaustive for the
residual and the audit contains no hidden derivation of out-of-scope
arithmetic.  The closure master theorem is therefore applicable to
\(\SelBSD\), and its residual conclusion is \(\rhoE=0\).
\end{proof}

\begin{theorem}[Composite signature with anti-tautology hardening]\label{thm:closure:composite-signature}
The composite recognition predicate attached to \(\SelBSD\) is
computable and falsifiable, carries a framework trace signature as a
fingerprint of the composite operator, and is anti-tautological: it
genuinely consumes all three recognition records.  In particular, the
three components are independent but not disposable.  Ablating
\(\GammaPD\), \(\GammaSP\), or \(\GammaGZ\) breaks the Strong-BSD
scalar factor package, so the composite does not collapse to a
tautological restatement of Strong BSD.\footnote{\raggedright Tracked in the
formalization harness as a conditional schema
\texttt{SixBirdsBSD.\allowbreak Closure.\allowbreak SelShell.\allowbreak compositeSignature}
-- a checked projection over an explicit proof-carrying record
(conditional on the imports and recognition sources), not an independent
Lean derivation. See App~\ref{app:formalization}.}
\end{theorem}

\begin{proof}
The composite predicate is a finite conjunction of the three readouts
introduced above, together with the shell data on which those readouts
are evaluated.  It is therefore computable in the framework sense: once
the shell and the recognition records are supplied, the predicate has a
definite list of components to inspect.  It is also falsifiable: failure
of any one component is a failure of the composite readout.

The substantive point is non-disposability.  Removing \(\GammaPD\)
leaves the local \(p\)-adic/Tamagawa factor undetermined; the shell no
longer knows that the descent comparison repairs the Tamagawa
coarsening by the required factor.  Removing \(\GammaSP\) leaves the
\(\Sha\)-persistence factor undetermined; the signed-Selmer determinant
comparison no longer supplies the component entering the arithmetic
product.  Removing \(\GammaGZ\) leaves the rank-\(r\)
regulator/fixity factor undetermined; there is then no assertion that
the analytic leading coefficient agrees with the normalized regulator
product.  Each recognition source governs a different scalar factor of
the residual, and each factor is required for the Strong-BSD scalar
package.  Hence the composite predicate depends genuinely on the
formed BSD layer's data and cannot be replaced by the bare sentence
``Strong BSD holds.''
\end{proof}

 \section{The Cassels--Tate comparison}
\label{sec:chi-ct-p}

The first supporting comparison used by the closure is the
Cassels--Tate Pfaffian comparison~\citep{Cassels1962,PoonenStoll1999}.  The import stack introduced in
Definition~\ref{def:closure:chi-ct-p-import} packages the relevant external
results as hypotheses: the Selmer-complex inputs of Sakamoto~\citep{Sakamoto2018}
and Macias Castillo--Sano~\citep{MaciasCastilloSano2026}, the determinant
and orientation inputs of Knudsen--Mumford and Nekov\'a\v{r}, and Flach's
pairing comparison.  Fisher--Schaefer--Stoll~\citep{FisherSchaeferStoll2010}
fix the classical 2-Selmer pairing normalization; they are not used as a
Fitting-ideal theorem.  This section states the
comparison in the form consumed by the closure.  It is conditional on
that stack; no part of the argument below reproves the imported
arithmetic geometry.

\begin{theorem}[$\chi_{CT,p}$ comparison (v5)]\label{thm:closure:chi-ct-p-comparison}
Let \(E/\Q\) be an elliptic curve such that \(\Sha(E)[p^\infty]\) is
finite, the Cassels--Tate pairing is nondegenerate, the standard
Sakamoto and Macias Castillo--Sano Selmer-complex hypotheses hold, and
\(p\nmid \Tam(E/\Q)\).  Conditional on the imported \(\chiCTp\)
structure, there is a canonical orientation-respecting Pfaffian
comparison map
\[
  \Pf_{\mathrm{Nek},p}^{\mathrm{or},\sqrt{\ }}:
  \sqrt{\Fitt}_{\mathrm{CT}}
    \bigl(H^1(C_p(E))/\mathrm{div},\,U_{2,2}\bigr)
  \longrightarrow
  \mathrm{vctp}_p
\]
from the square-root Cassels--Tate Fitting ideal on the quotient by the
divisible subgroup, paired by the Nekov\'a\v{r} cup-product
\(U_{2,2}\) identified with the Flach pairing, to the
Cassels--Tate Pfaffian scalar line \(\mathrm{vctp}_p\).  This map sends
the square-root Stark generator~\citep{MazurRubin2004} \(\sqrt{F}_E^p\) to the
Cassels--Tate height \(h_p^{\mathrm{CT}}(E)\), and the Nekov\'a\v{r}
sign convention satisfies \(\Sigma_{\mathrm{NekCT}}=+1\).\footnote{Tracked in the formalization harness as a conditional schema
\texttt{SixBirdsBSD.\allowbreak Closure.\allowbreak ChiCTp.\allowbreak chiCTpComparison}
-- a checked projection over an explicit proof-carrying record
(conditional on the imported \(\chiCTp\) stack), not an independent Lean
derivation. See App~\ref{app:formalization}.}
\end{theorem}

\begin{proof}
Assume the imported \(\chiCTp\) structure.  Its named components supply
the ingredients in the statement.  The Selmer-complex hypotheses and
the Cassels--Tate nondegeneracy input place the quotient
\(H^1(C_p(E))/\mathrm{div}\) in the self-dual setting in which a
Pfaffian comparison is defined.  The determinant-line and orientation
inputs identify the square-root Fitting ideal
\(\sqrt{\Fitt}_{\mathrm{CT}}(H^1(C_p(E))/\mathrm{div},U_{2,2})\) as the
source of the comparison map.  Here the square-root Fitting ideal is
the square-root determinant object attached to the Cassels--Tate
pairing; the notation \(U_{2,2}\) records the Nekov\'a\v{r}
cup-product, and the imported Flach comparison identifies it with the
classical Cassels--Tate pairing used on the target side.

The Pfaffian input in the stack then gives the canonical
orientation-respecting map
\[
  \Pf_{\mathrm{Nek},p}^{\mathrm{or},\sqrt{\ }}:
  \sqrt{\Fitt}_{\mathrm{CT}}
    \bigl(H^1(C_p(E))/\mathrm{div},U_{2,2}\bigr)
  \longrightarrow
  \mathrm{vctp}_p .
\]
The square-root Stark-generator identification supplies the image
formula: the distinguished generator \(\sqrt{F}_E^p\) maps to the
Cassels--Tate height \(h_p^{\mathrm{CT}}(E)\).  Finally, the sign
closure component of the same imported record fixes the orientation
sign as \(\Sigma_{\mathrm{NekCT}}=+1\).  Combining these imported
components proves exactly the displayed comparison.  The proof is
therefore an assembly under the stated hypotheses; it is precisely as
strong as the imported theorem stack and does not rederive the
Sakamoto, Macias Castillo--Sano, Nekov\'a\v{r}, Flach, or classical
Cassels--Tate comparison inputs.
\end{proof}

 \section{The OC eta-formula}
\label{sec:eta-formula}

This section records the second supporting conditional theorem: an
overconvergent \(\eta\)-formula at additive odd primes.  It is a
\(p\)-adic \(L\)-value identity used by the closure in the shallow
additive odd-prime regime.  Its inputs are separate from the
Cassels--Tate comparison of the previous section; they are collected in
their own published-import structure.

\begin{definition}[OC $\eta$ published imports]\label{def:closure:eta-published-imports}
The published-import structure for the overconvergent \(\eta\)-formula
keeps the following inputs separate: Bella\"iche's critical-slope
construction and moment formula~\citep{Bellaiche2012}; the Pollack--Stevens overconvergent modular-symbol theory~\citep{PollackStevens2011,ColemanMazur1998}; the
critical \(p\)-adic \(L\)-function input~\citep{BenoisBuyukboduk2024}; Lang--Wake~\citep{LangWake2025}, as a
distinct import from the critical-\(p\)-adic-\(L\)-function work;
Lorenzini's component-group input~\citep{Lorenzini1993}; Ling's refinement~\citep{Ling1997} of the local
component calculation; and one cascade-internal nonvanishing and
primitivity conjecture for the cuspidal-layer evaluation.\footnote{Encoded in Lean as a typed
structure carrier
\texttt{SixBirdsBSD.\allowbreak Closure.\allowbreak EtaFormula.\allowbreak etaPublishedImports}
bundling the imported results as hypothesis fields; not a Lean axiom.
See App~\ref{app:formalization}.}
\end{definition}

\begin{theorem}[Conditional OC $\eta$-formula]\label{thm:closure:oc-eta-formula}
Let \(E/\Q\) lie in the shallow additive odd-prime regime, including
the prime-square anchors and the verified multi-prime extension, such
as \(36a1\) at \(p=3\).  Let \(p\) be an odd additive prime.  Let
\(\beta\) be the nonzero \(U_p\)-eigenvalue of the chosen additive
\(p^2\)-stabilization, and assume the local relation \(\beta^2=-p\).  Let
\(T_{\mathrm{leading\text{-}Bellaiche}}(E,p)\) be the Bella\"iche
leading term of the overconvergent expansion, and let
\(\mathrm{ord}_T(E,p)\) be its order in the expansion parameter.
Conditional on Definition~\ref{def:closure:eta-published-imports},
\[
  \etap(E)
  =
  \bigl(p\,|E(\Q)_{\tors}|\bigr)^{\mathrm{ord}_T(E,p)-1}
  \,\beta^{-2}\,
  T_{\mathrm{leading\text{-}Bellaiche}}(E,p).
\]
The anchor evidence is anchor-level evidence only: the OC
\(\eta\)-formula sub-residual is verified at \(49a1\) at \(p=7\)
with \(\eta=-2\), at \(36a1\) at \(p=3\) with \(\eta=+2\), and at
\(121b1\) at \(p=11\) with \(\eta=-1\), not as a uniform
family-level statement.\footnote{\raggedright Tracked in the formalization
harness as a conditional schema
\texttt{SixBirdsBSD.\allowbreak Closure.\allowbreak EtaFormula.\allowbreak ocEtaFormula}
-- a checked projection over an explicit proof-carrying record
(conditional on the published imports), not an independent Lean
derivation. See App~\ref{app:formalization}.}
\end{theorem}

\begin{proof}
Assume the published-import structure of
Definition~\ref{def:closure:eta-published-imports}.  The chosen additive
stabilization fixes \(\beta\) with \(\beta^2=-p\); Bella\"iche's moment
construction contributes the \(\beta^{-2}\) factor at level \(p^2\).
The overconvergent modular-symbol input of Pollack--Stevens supplies the overconvergent expansion in which the
Bella\"iche leading term \(T_{\mathrm{leading\text{-}Bellaiche}}(E,p)\)
and the order \(\mathrm{ord}_T(E,p)\) are defined.  The critical
\(p\)-adic \(L\)-function input identifies this leading term with the
corresponding \(p\)-adic \(L\)-value contribution.

The remaining published inputs, namely Lang--Wake (2025), Lorenzini,
and Ling, together with the single cascade-internal nonvanishing and
primitivity conjecture, supply the normalization of the local
component and torsion factors.  Combining those normalizations with the
critical-slope relation gives exactly
\[
  \etap(E)
  =
  \bigl(p\,|E(\Q)_{\tors}|\bigr)^{\mathrm{ord}_T(E,p)-1}
  \,\beta^{-2}\,
  T_{\mathrm{leading\text{-}Bellaiche}}(E,p).
\]
The three numerical anchors are finite checks of this displayed
identity at the stated curves and primes.  They are not used as a
uniform proof over the shallow additive odd-prime regime.  Thus the
formula is exactly as strong as the published-import stack together
with the named nonvanishing/primitivity conjecture; it does not
rederive any of the imported results.
\end{proof}

 \section{The rank-one cascade import}
\label{sec:t-cascade}

This section records the conditional rank-\(\leq 1\) input used by the
closure: a \(p\)-part Strong-BSD valuation identity in ranks zero and
one.  The result is conditional on a large named external stack of
Iwasawa-theoretic, Heegner-point, and Selmer-complex results.  The
stack has two branches, one for CM curves and one for non-CM curves.

\begin{definition}[$T_{\mathrm{CASCADE}}$ import stack]\label{def:closure:t-cascade-import}
The cascade import is a proof-carrying structure bundling a stack of
named external results.  In the CM branch it packages the standard CM
Iwasawa-theoretic inputs~\citep{Rubin1991}, including the ordinary-split and
supersingular-inert main-conjecture inputs, together with the
archimedean, Heegner, and uniformization ingredients needed to compare
the \(p\)-adic and classical sides.  In the non-CM branch it packages
the ordinary and Euler-system inputs of Skinner--Urban~\citep{SkinnerUrban2014}, Wei Zhang~\citep{WeiZhang2014},
Burungale--Castella--Skinner--Tian~\citep{BurungaleCastellaSkinnerTian2022},
CHKLL~\citep{CHKLL2025}, and Yan--Zhu~\citep{YanZhu2026}.  The
closure takes this stack as a hypothesis.\footnote{Encoded in Lean as a typed
structure carrier
\texttt{SixBirdsBSD.\allowbreak Closure.\allowbreak TCascade.\allowbreak tCascadeImport}
bundling the imported results as hypothesis fields; not a Lean axiom.
See App~\ref{app:formalization}.}
\end{definition}

\begin{theorem}[Conditional $p$-part SBSD, rank $\leq 1$]\label{thm:closure:t-cascade-rank-le-1}
Let \(E/\Q\) have analytic rank \(r\in\{0,1\}\).  In the CM branch,
assume that \(E\) has complex multiplication by \(K\), and that either
\(p\) is a good ordinary prime split in \(K\), or \(p\) is a good
supersingular prime inert in \(K\) with \(p\geq 5\).  Assume also that
\(\Sha(E)[p^\infty]\) is finite, the Cassels--Tate pairing is
nondegenerate, the Sakamoto and Macias Castillo--Sano Selmer-complex hypotheses
hold, \(p\nmid\Tam(E/\Q)\), the standard ramification and CM
Galois-image hypotheses hold, and a Heegner point is supplied when
\(r=1\).

In the non-CM branch, assume that \(E\) is non-CM, \(r\in\{0,1\}\),
the representation \(\rho_{E,p}\) is surjective~\citep{Serre1972}, \(p\) is good
ordinary with \(p\nmid N\), and \(p\nmid\Tam(E/\Q)\).  Assume the
standard Skinner--Urban ordinary-local, residual-irreducibility,
ramification, and level hypotheses; assume \(\Sha(E)[p^\infty]\) is
finite, the Cassels--Tate pairing is nondegenerate, and the
Sakamoto and Macias Castillo--Sano hypotheses hold.  When \(r=1\), also assume a
Heegner field satisfying the Gross--Zagier--Kolyvagin hypotheses.

Then, conditional on Definition~\ref{def:closure:t-cascade-import}, the
\(p\)-part valuation identity
\[
  v_p\!\left(\frac{L^{(r)}(E,1)}{r!\,\Omega\,R}\right)
  =
  v_p\!\left(\frac{\#\Sha(E/\Q)\,\Tam(E)}
    {|E(\Q)_{\tors}|^2}\right)
\]
holds, where \(R\) is the N\'eron--Tate regulator.  Thus the cascade
supplies the \(p\)-part of the Strong-BSD identity in ranks zero and
one.\footnote{\raggedright Tracked in the formalization harness as a
conditional schema
\texttt{SixBirdsBSD.\allowbreak Closure.\allowbreak TCascade.\allowbreak tCascadeRankLeOne}
-- a checked projection over an explicit proof-carrying record
(conditional on the imported cascade stack), not an independent Lean
derivation. See App~\ref{app:formalization}.}
\end{theorem}

\begin{proof}
Assume the cascade import stack.  In the CM branch, the CM
Iwasawa-theoretic main-conjecture inputs identify the \(p\)-adic
\(L\)-function side with the corresponding Selmer-theoretic side in
the ordinary-split and supersingular-inert cases.  The Selmer-complex
and Cassels--Tate hypotheses place the \(p\)-primary part of
\(\Sha\) in the required nondegenerate comparison setting.  When
\(r=1\), the Heegner-point input supplies the regulator term; in rank
zero no Heegner regulator input is needed.  The standard ramification
and CM Galois-image hypotheses are exactly the local hypotheses under
which the imported CM comparison applies.

In the non-CM branch, the Skinner--Urban input gives one divisibility
between the \(p\)-adic \(L\)-value side and the Selmer side.  The
opposite divisibility is supplied by the Euler-system and
Gross--Zagier--Kolyvagin inputs, with the Gross--Zagier--Kolyvagin
Heegner field used only in rank one.  In rank zero, the argument uses
the rank-zero part of the same imported ordinary and Euler-system
comparison and does not assume a rank-one Heegner field.

In both branches, the hypotheses \(p\nmid\Tam(E/\Q)\), finiteness of
\(\Sha(E)[p^\infty]\), and nondegeneracy of the Cassels--Tate pairing
make the local and \(\Sha\)-terms comparable by \(p\)-adic units in the
imported comparison.  After these unit ambiguities are removed, the two
divisibilities in each branch pin down exactly the displayed valuation
identity
\[
  v_p\!\left(\frac{L^{(r)}(E,1)}{r!\,\Omega\,R}\right)
  =
  v_p\!\left(\frac{\#\Sha(E/\Q)\,\Tam(E)}
    {|E(\Q)_{\tors}|^2}\right).
\]
The proof is therefore an assembly of the imported cascade stack under
the stated branch hypotheses.  It is precisely as strong as that stack:
it does not rederive Skinner--Urban, Wei Zhang, the CM main conjecture,
or Gross--Zagier--Kolyvagin.
\end{proof}

 \section{Typed obstructions}
\label{sec:obstructions}

The closure isolates three recognition sources because the
theorem-grade literature does not presently supply every identity that
would be needed for an unconditional Strong-BSD proof.  We call that
literature ``Mode A'': the peer-reviewed arithmetic-geometry literature
audited as external provenance, in the same sense as
\citet{TsiokosBSDApparatus2026}.  The two statements
below are audit results about the current reach of that literature.
They are not non-existence claims.  The missing content may well exist;
the point is that the audited literature does not currently provide it
in theorem-grade form.  The closure therefore carries the missing
arithmetic content as recognition sources rather than deriving it.

\begin{theorem}[Higher-rank Gross--Zagier audit no-go]\label{thm:closure:t-e12-no-go}
No theorem-grade peer-reviewed identity of the form
\[
  \frac{L^{(r)}(E,1)}{r!}
  =
  \kapr(E)\,\det\langle Z_i,Z_j\rangle_{\mathrm{NT}}\,
  \Omega_E\,c_E
\]
is available in the audited Mode A literature for non-CM \(E/\Q\) of
rank \(r\geq 2\), with explicit algebraic cycles \(Z_i\) and closed
forms for \(\kapr(E)\) and \(c_E\).  This is the missing rank-\(r\)
generalized Gross--Zagier identity.  The audited Mode A approaches
converge on this gap.  This is explicitly not a non-existence claim:
such an identity may exist, but it is not presently supplied by the
audited theorem-grade literature.
\end{theorem}

\begin{theorem}[Additive small-prime main-conjecture audit no-go]\label{thm:closure:t-bad-no-go}
No theorem-grade direct additive-prime Iwasawa main conjecture is
available in the audited Mode A literature for \(E/\Q\) additive at
\(p\in\{2,3\}\), with explicit local correction factors and
characteristic-ideal equality.  In the audit this gap separates into a
small list of local residual checks together with a Kodaira-type
admissibility screen.  Again, this is an audit-result no-go, not a
claim that such a theorem cannot exist.
\end{theorem}

\begin{remark}[Binding residuals carried by the recognition sources]\label{rmk:closure:mode-b-residuals}
The cascade also produces three internal residuals: the rank-\(r\)
regulator/fixity residual, the signed-Selmer corestriction residual,
and the \(p\)-adic \(L\)-norm identity residual.  These are not
independent theorems of the paper.  They are the reorganized arithmetic
content carried by the three recognition sources of
\Cref{sec:recognition-sources}: the higher-rank fixity source carries
the regulator/fixity residual, the signed-Selmer persistence source
carries the corestriction residual, and the \(p\)-adic descent source
carries the \(p\)-adic norm identity residual.  Thus they are recorded
as support-only typed-open content, not as additional derived claims.
\end{remark}

 \section{The conditional Strong-BSD landing}
\label{sec:landing}

We now assemble the closure.  The preceding sections supplied the
trace shell \(\SelBSD\), the readout-level equivalence between
\(\rhoE=0\) and the Strong-BSD scalar identity, the applicability of
the closure master theorem, the non-disposability of the three
recognition records, the three recognition sources themselves, and the
headline imported theorem stacks.  The result below is the promised conditional
landing.
\Cref{tab:landing-chain,fig:landing-chain-diagram} give the chain before the
formal statement.
\begin{table}[tbp]
\centering
\small
\caption{The conditional landing chain in companion-form prose.}
\label{tab:landing-chain}
\begin{tabularx}{\textwidth}{@{}>{\raggedright\arraybackslash}p{0.11\textwidth}
>{\raggedright\arraybackslash}p{0.35\textwidth}
>{\raggedright\arraybackslash}X@{}}
\toprule
Step & Premise & Conclusion \\
\midrule
1 & The standing shell-closure record, the three recognition sources, and
the headline imported theorem stacks are supplied to \(\SelBSD\). & The
shell has named records for every scalar factor of \(\rhoE\). \\
\addlinespace
2 & The applicability and no-smuggling audits pass. & The closure master
theorem applies to the shell. \\
\addlinespace
3 & The master theorem applies to \(\SelBSD\). & The residual readout
vanishes: \(\rhoE=0\). \\
\addlinespace
4 & The readout-level equivalence identifies \(\rhoE=0\) with the displayed
scalar identity. & Strong BSD follows under the stated hypotheses. \\
\bottomrule
\end{tabularx}
\end{table}
 \begin{figure}[tbp]
\centering
\begin{tikzpicture}[
  font=\small,
  input/.style={draw, rounded corners=2pt, align=center, minimum width=4.7cm,
    minimum height=1.35cm, fill=black!3},
  box/.style={draw, rounded corners=2pt, align=center, minimum width=3.0cm,
    minimum height=1.0cm, fill=black!3},
  result/.style={draw, rounded corners=2pt, align=center, minimum width=4.0cm,
    minimum height=1.35cm, fill=black!8},
  arrow/.style={-{Latex[length=2mm]}, thick}
]
\node[input] (inputs) at (-4.9,0.25) {shell-closure record\\
  three recognition records\\headline imports};
\node[box] (master) at (0,0.25) {closure master\\theorem applies};
\node[result] (landing) at (4.9,0.25) {\(\rhoE=0\)\\
  {\scriptsize readout equivalence}\\Strong-BSD identity};

\draw[arrow] (inputs) -- (master);
\draw[arrow] (master) -- (landing);
\node[align=center, font=\footnotesize] at (0,-1.15)
  {conditional on the supplied records, the shell residual vanishes and the scalar identity follows};
\end{tikzpicture}
\caption{The conditional landing chain from the shell-closure record,
recognition records, and headline imports to the Strong-BSD scalar identity.}
\label{fig:landing-chain-diagram}
\end{figure}

\begin{theorem}[Conditional Strong BSD]\label{thm:closure:strong-bsd-conditional}
Assume the standing closure assumption for \(\SelBSD\): the formed
object carries its admissibility, anti-tautology, and lawfulness data.
Assume the three recognition sources of
Definitions~\ref{def:closure:gamma-padic-descent},
\ref{def:closure:gamma-sha-persistence}, and
\ref{def:closure:gamma-higher-gz-fixity}.  Assume also the three
headline import structures of Definitions~\ref{def:closure:chi-ct-p-import},
\ref{def:closure:a-e-import}, and
\ref{def:closure:beilinson-import}.  Then the Strong-BSD scalar
identity holds for \(E/\Q\):
\[
  \frac{L^{(r)}(E,1)}{r!}
  =
  \frac{\#\Sha(E/\Q)\,\Reg_{\mathrm{NT}}(E)\,\Omega_E
    \prod_v c_v(E)}
       {|E(\Q)_{\tors}|^2}.
\]
Outside the framework, this is exactly the conditional theorem
\[
  \begin{aligned}
  &\SelBSD\text{-closure}
    \wedge\ \GammaPD \wedge \GammaSP \wedge \GammaGZ\\
  &\quad{}\wedge\ \chiCTp\text{-imports}
    \wedge\ \AofE\text{-imports}\\
  &\quad{}\wedge\ \text{Beilinson-rank-}{\geq}2\text{-imports}\\
  &\Longrightarrow \mathrm{Strong\ BSD},
  \end{aligned}
\]
not an unconditional proof of BSD.\footnote{\raggedright Tracked in the
formalization harness as a conditional schema
\texttt{SixBirdsBSD.\allowbreak Closure.\allowbreak Landing.\allowbreak strongBSDConditional}
-- a checked projection over an explicit proof-carrying record whose
hypotheses are the recognition sources and headline imports (a genuine
implication chain through the shell-readout theorems), not an
independent Lean derivation. See App~\ref{app:formalization}.}
\end{theorem}

\begin{proof}
We follow the shell-readout chain.  By
\Cref{thm:closure:master-theorem-applicability}, the three recognition
sources are correctly instantiated on \(\SelBSD\), exhaust the scalar
factors of the residual \(\rhoE\), and pass the no-smuggling audit.
Thus the closure master theorem applies to the shell and gives
\[
  \rhoE=0.
\]
Let us spell out the factor accounting.  The higher-rank fixity source
\(\GammaGZ\) supplies the assertion \(\psi_-(E,r)=0\).  By the
\(\kapr\) normalization of \citet{TsiokosBSDApparatus2026},
\[
  \kapr(E)=\frac{\#\Sha(E/\Q)}{|E(\Q)_{\tors}|^2},
\]
so the product
\(\kapr(E)\,\Reg_{\mathrm{NT}}(E)\,\Omega_E\,\Tam(E)\) is the
rank-\(r\) arithmetic part of \(\rhoE\), with \(\Tam(E)=\prod_v c_v(E)\).
The \(p\)-adic descent source \(\GammaPD\) supplies the local
\(p\)-adic/Tamagawa factor; the Cassels--Tate and Selmer-complex
imports support the comparison of that local factor with the
Cassels--Tate/Selmer-complex side.  The signed-Selmer persistence
source \(\GammaSP\) supplies the \(\Sha\)-persistence factor.  The
Beilinson import supplies the rank-\(\geq 2\) archimedean
determinant/\(L\)-value identification.  Together these are exactly
the scalar factors of \(\rhoE\), so the master-theorem conclusion is
the vanishing of the residual.

Now apply the readout-level equivalence
\Cref{thm:closure:pi-bsd-iff-strong-bsd}.  Since \(\rhoE\) is defined as
\[
  \frac{L^{(r)}(E,1)}{r!}
  -
  \frac{\#\Sha(E/\Q)\,\Reg_{\mathrm{NT}}(E)\,\Omega_E
    \prod_v c_v(E)}
       {|E(\Q)_{\tors}|^2},
\]
the equality \(\rhoE=0\) is precisely the displayed Strong-BSD scalar
identity.  Hence the identity follows.

Finally, \Cref{thm:closure:composite-signature} records the
anti-tautology check: the composite readout genuinely consumes all
three recognition sources.  Removing \(\GammaPD\) leaves the local
\(p\)-adic/Tamagawa factor undetermined; removing \(\GammaSP\) leaves
the \(\Sha\)-persistence factor undetermined; removing \(\GammaGZ\)
leaves the regulator/fixity factor undetermined.  Thus the conclusion
is not a tautological restatement of Strong BSD.  Under the stated
closure assumption, recognition sources, and imported stacks, the
Strong-BSD scalar identity holds.
\end{proof}

\begin{nonclaim}\label{nonclaim:closure:landing-conditional}
This is not an unconditional classical proof of BSD, scalar or strong.
The result is theorem-grade within the Six Birds closure discipline,
conditional on the standing shell-closure record, the named recognition
sources, and the headline imported theorem stacks; outside the framework it
is the conditional implication
displayed in the theorem above.  The rank-\(\geq 2\) generalized
Gross--Zagier identity and the additive \(p\leq 3\) Iwasawa main
conjecture remain genuinely open in the literature.
\end{nonclaim}

 \section{BSD as a non-eliminative AOR instance}
\label{sec:aor-instance}

The conditional landing of \Cref{sec:landing} gives the scalar
Strong-BSD identity under the standing shell-closure record, the named
recognition sources, and the headline imports.
This final technical section records the corresponding
Audited Operational Realisability reading.  In the sense of
\cite{TsiokosAOR2026}, an AOR instance is a refinement-stable,
audit-closed membership predicate: informally, the framework's
``holds in reality'' assertion for an object, meaning that the
residual defects left by refinement have all been accounted for by
approved records.  The reading here is non-eliminative.  The open
arithmetic content is not erased; it is recorded as recognition-source
content, while the mechanical pieces and imported theorem stacks are
accounted for by construction or by approved external citation.

\begin{definition}[BSD AOR-instance object]\label{def:closure:aor-sel-instance}
The AOR-instance object attached to \(\SelBSD\) is the carrier whose
membership witness is the bundle of closure records for the BSD trace
shell.  The bundle has two kinds of entries.  The mechanical entries
are the construction of \(\SelBSD\), the \(\kapr\) normalization, and
the imported theorem stacks recorded through an approved
external-record citation chain.  The substantive entries are the three
recognition-source records of \Cref{sec:recognition-sources}.  The
object is built using the AOR primitives of \cite{TsiokosAOR2026}:
refinement-stable membership, structural-defect accounting, and
discharge records.
\end{definition}

\begin{theorem}[Mechanical-record discharges]\label{thm:closure:aor-mechanical-records}
The mechanical membership components discharge cleanly.  The
\(\SelBSD\) construction and the \(\kapr\) normalization are
discharged by construction.  The \(\chiCTp\), \(\AofE\), and
Beilinson import structures are discharged by an approved
external-record citation chain: they are cited as imported theorem
stacks, not reproved.

For the normalization record, we use \citet{TsiokosBSDApparatus2026}.
Fix the cascade convention
\[
  c_E := \Tam(E)=\prod_{\ell} c_{\ell}(E).
\]
Then the cascade leading-term form
\[
  \frac{L^{(r)}(E,1)}{r!}
  =
  \kapr(E)\,\Omega_E\,\Reg_r(E)\,c_E
\]
agrees with the standard Strong-BSD leading-coefficient form
\[
  \frac{L^{(r)}(E,1)}{r!}
  =
  \Omega_E\,\Reg_r(E)\,
  \frac{\#\Sha(E/\Q)\,\Tam(E)}{\#E(\Q)_{\tors}^{\,2}}
\]
if and only if
\[
  \kapr(E)=
  \frac{\#\Sha(E/\Q)}{\#E(\Q)_{\tors}^{\,2}}.
\]
This is a restatement in cascade coordinates, not a proof of the
Strong-BSD identity.
\footnote{Verified in Lean as
\texttt{SixBirdsBSD.Closure.AORInstance.aorMechanicalRecords};
see App~\ref{app:formalization}.}
\end{theorem}

\begin{proof}
The shell itself is one of the constructed objects of the closure
architecture, so its mechanical entry is accounted for by construction.
The displayed \(\kapr\) equivalence is likewise a constructed
normalization record: once the cascade convention \(c_E=\Tam(E)\) is
fixed, it identifies the scalar that makes the cascade leading-term
form and the standard Strong-BSD leading-coefficient form the same
formula.  It does not assert that the formula is true independently of
the conditional landing.

It remains only to account for the imported theorem stacks.  The
\(\chiCTp\) Cassels--Tate stack, the \(\AofE\) Selmer-complex pairing
stack, and the Beilinson comparison stack enter as named external
records.  They are therefore discharged by an approved citation chain,
not by a derivation inside the closure.  These three citation records,
together with the constructed shell and the \(\kapr\) normalization,
exhaust the mechanical part of the AOR membership witness.
\end{proof}

\begin{theorem}[Recognition-source discharges]\label{thm:closure:aor-recognition-discharge}
The three recognition sources discharge as bridged records.  Here a
bridged record means a recognition residual that is reclassified as
supplied closure content rather than eliminated by a derivation.  Each
\(\Gamma_{\mathrm{BSD}}^{*}\) is a primary structural defect of
source type, carrying its forced secondary defects of role, target,
transport, and limit, and is discharged by such a bridged record.  In
particular, none of the three recognition sources is eliminated.
\footnote{Verified in Lean as
\texttt{SixBirdsBSD.Closure.AORInstance.aorRecognitionDischarge};
see App~\ref{app:formalization}.}
\end{theorem}

\begin{proof}
Each of the three recognition sources is an open arithmetic obligation
that the closure names and consumes but does not derive.  In the audit
presentation, such an obligation is therefore recorded as a
source-type structural defect, together with the secondary information
that specifies its role in the argument, its target factor in the BSD
residual, its transport through the shell, and the limit of what the
closure claims about it.

The discharge is by bridging.  That is, the record says that the
obligation is supplied as recognition content and is available to the
conditional landing; it does not say that the obligation has been
proved away.  This applies to the \(p\)-adic descent source, the
signed-Selmer persistence source, and the higher-rank fixity source.
Thus all three recognition records are accounted for, and all three
remain non-eliminated.
\end{proof}

\Cref{tab:aor-discharge-register} lists the mechanical, imported, and
recognition-source atoms that enter the local AOR membership witness.
\begin{table}[tbp]
\centering
\small
\caption{The eight AOR discharge atoms for the BSD trace shell.  The
recognition-source rows are bridged, not zeroed.}
\label{tab:aor-discharge-register}
\begin{tabularx}{\textwidth}{@{}>{\raggedright\arraybackslash}p{0.31\textwidth}
>{\raggedright\arraybackslash}p{0.24\textwidth}
>{\raggedright\arraybackslash}X@{}}
\toprule
Residual atom & Primary type & Status \\
\midrule
\(\SelBSD\) construction & mechanical & by construction \\
\(\kapr\) normalization & mechanical & by construction, citing the
apparatus paper \\
\(\chiCTp\) import stack & imported theorem stack & approved external
record \\
\(\AofE\) Selmer-complex substrate & imported theorem stack & approved
external record \\
Beilinson rank-\(\geq2\) comparison & imported theorem stack & approved
external record \\
\(\GammaPD\) & recognition source & bridged, not eliminated \\
\(\GammaSP\) & recognition source & bridged, not eliminated \\
\(\GammaGZ\) & recognition source & bridged, not eliminated \\
\bottomrule
\end{tabularx}
\end{table}

\begin{theorem}[AOR-instance membership]\label{thm:closure:aor-instance}
\(\SelBSD\) is refinement-stable audit-closed, hence an AOR instance,
under the declared presentation: its declared audit records discharge
the structural defects that remain under refinement.  This is the
framework's ``Strong BSD holds in reality'' statement, and it is
non-eliminative.  The recognition records are reclassified as bridged
records, not discharged of their open arithmetic.
\footnote{Verified in Lean as
\texttt{SixBirdsBSD.Closure.AORInstance.aorInstance} over a narrowed
representation -- the local syntactic refinement-stable predicate, not
the full Audited Operational Realisability meta-theory; the narrowing
is recorded in App~\ref{app:formalization}.}
\end{theorem}

\begin{proof}
By \Cref{thm:closure:aor-mechanical-records}, the mechanical
components of the BSD trace shell are accounted for: the shell
construction and the \(\kapr\) normalization are discharged by
construction, while the imported theorem stacks are discharged by
approved external citation.  By
\Cref{thm:closure:aor-recognition-discharge}, the three recognition
sources are accounted for as bridged records.

These two groups exhaust the declared structural defects that persist
under refinement of the BSD trace shell.  The mechanical records are
closed by construction or citation; the recognition records are
carried forward as supplied recognition content.  Therefore every
remaining defect is matched by a declared audit record, so the shell is
refinement-stable audit-closed under the declared presentation.

The last clause is essential.  The membership statement does not
convert the recognition sources into classical proofs of the open
arithmetic obligations.  It records that those obligations are present,
named, and bridged in the audit, which is exactly the non-eliminative
AOR reading.
\end{proof}

\begin{remark}[AOR partial status]\label{rmk:closure:aor-partial-status}
The AOR-instance reading does not close the open obligations.  The
rank-\(\geq 2\) generalized Gross--Zagier identity and the additive
\(p\leq 3\) Iwasawa main conjecture remain open in the literature.
What the AOR reading records is that the BSD carrier is audit-closed
under the declared refinement-stable presentation, with the
recognition records reclassified as bridged content rather than
discarded.  Only the local syntactic refinement-stable predicate is
mechanized here, not the full AOR meta-theory.
\end{remark}
 \section{Scope and non-claims}\label{sec:scope}


The closure developed in this paper is deliberately conditional.  The
point of the construction is to isolate exactly which recognition
sources and imported theorem stacks would close the Strong-BSD scalar
identity inside the Six Birds closure discipline, and to record what
is not being claimed.  The following seven non-claims are part of the
statement of the paper.  They fix the intended reading of the theorem,
the role of the recognition sources, the status of the imports, and
the meaning of the final AOR membership assertion.

\begin{itemize}[leftmargin=*]
\item[\textbf{NC-1.}]
\textbf{Not unconditional BSD.}  This is not an unconditional
classical proof of BSD, scalar or strong.  The result is theorem-grade
within the Six Birds closure discipline, conditional on the standing
shell-closure record, the named recognition sources, and the headline
imported theorem stacks.  Outside the framework it reads as the
conditional theorem
\[
  \begin{aligned}
  &\SelBSD\text{-closure}
    \wedge\ \GammaPD \wedge \GammaSP \wedge \GammaGZ\\
  &\quad{}\wedge\ \chiCTp\text{-imports}
    \wedge\ \AofE\text{-imports}
    \wedge\ \text{Beilinson-imports}\\
  &\Longrightarrow \mathrm{Strong\ BSD}.
  \end{aligned}
\]
This is the headline reading of the paper.  The closure is a precise
conditional implication, not a replacement for the missing
unconditional arithmetic theorems.

\item[\textbf{NC-2.}]
\textbf{Recognition sources not derived.}  No recognition source is
derived from framework primitives.  Each \(\Gamma_{\mathrm{BSD}}^{*}\)
is closure content of the formed BSD layer: a typed carrier whose
mathematical payload is supplied as a hypothesis and then consumed by
the shell.  The formalization records these carriers as explicit
structure data, not as axioms and not as hidden constants.  Thus the
three sources are neither proved by the framework nor smuggled into
the conclusion; they are named, typed, and kept visible throughout the
argument.

\item[\textbf{NC-3.}]
\textbf{Rank-\(\geq 2\) Gross--Zagier remains open.}  We do not prove
the rank-\(\geq 2\) generalized Gross--Zagier identity.  The relevant
missing identity is the higher-rank comparison between the leading
\(L\)-value and a Neron--Tate regulator determinant of algebraic
cycles, with the necessary normalizing factors made explicit.  The
no-go result used in this paper is an audit result: in the audited
theorem-grade literature, the needed number-field identity is not
available in the required form.  This is not a non-existence claim.
It says only that the present literature does not supply the theorem
needed to remove the corresponding recognition source.

\item[\textbf{NC-4.}]
\textbf{Additive \(p\leq 3\) Iwasawa main conjecture remains open.}
We do not prove the additive \(p\leq 3\) Iwasawa main conjecture with
the explicit local correction factors and characteristic-ideal
identity required by the closure.  This is again an audit-result
no-go, not a claim that such a theorem cannot exist.  The paper treats
the missing small-additive-prime input as open arithmetic content and
does not derive it from the closure framework.  In particular, the
odd-additive recognition carriers do not by themselves discharge the
\(p=2\) part of this obstruction; any such discharge would have to be a
separate supplied record or a future imported theorem.

\item[\textbf{NC-5.}]
\textbf{Imports not eliminated.}  The \(\chiCTp\) comparison remains
conditional on its imported Cassels--Tate stack.  Likewise, the
\(\AofE\) Selmer-complex pairing substrate and the Beilinson
rank-\(\geq 2\) comparison enter as imported theorem stacks.  These
structures are sourced assumptions for this paper, not results proved
here.  The closure derives consequences from them; it does not
eliminate them, reconstruct them, or claim independent proofs of the
classical arithmetic geometry they package.

\item[\textbf{NC-6.}]
\textbf{Anchor-level evidence is not family-level.}  The
overconvergent \(\eta\)-formula sub-residual is verified at the
anchors \(49\mathrm{a}1@7\), \(36\mathrm{a}1@3\), and
\(121\mathrm{b}1@11\).  Those checks are anchor-level numerical
evidence.  They are not a uniform family-level statement over all
elliptic curves or all additive odd primes in the relevant regime.
The conditional formula in \Cref{sec:eta-formula} remains conditional
on its published-import stack and its named nonvanishing/primitivity
input.

\item[\textbf{NC-7.}]
\textbf{AOR is non-eliminative.}  The AOR-instance reading does not
close the open obligations.  It reclassifies the closure records: the
recognition discharges are bridged records, meaning the corresponding
open arithmetic content is carried as supplied recognition content
rather than eliminated.  The rank-\(\geq 2\) generalized
Gross--Zagier identity and the additive \(p\leq 3\) Iwasawa main
conjecture therefore remain open after the AOR reading.  Moreover,
only the local syntactic refinement-stable predicate is mechanized in
this paper, not the full Audited Operational Realisability
meta-theory.
\end{itemize}

These non-claims also govern the wording of the paper.  We do not say
that the paper proves BSD, proves the additive-prime Iwasawa main
conjecture, or proves the higher-rank Gross--Zagier identity.  The
paper lands Strong BSD as a conditional closure.  The standing shell-closure
record is assumed; recognition sources are supplied and carried as
hypotheses; headline imported theorem stacks are imported and cited; the
AOR reading records a non-eliminative
audit-closed membership statement under those declared records.
\begin{table}[tbp]
\centering
\footnotesize
\caption{Closure-paper nonclaims.  These are wording boundaries for the
headline theorem, recognition sources, imports, and AOR reading.}
\label{tab:nonclaims}
\begin{tabularx}{\textwidth}{@{}>{\raggedright\arraybackslash}p{0.12\textwidth}
>{\raggedright\arraybackslash}p{0.25\textwidth}
>{\raggedright\arraybackslash}X@{}}
\toprule
Code & Boundary & Short reading \\
\midrule
NC-1 & Not unconditional BSD & The theorem is conditional on the standing
shell-closure record, named recognition sources, and headline imports. \\
NC-2 & Recognition sources not derived & The \(\Gamma_{\mathrm{BSD}}^*\)
carriers are supplied hypotheses, not framework consequences. \\
NC-3 & Rank-\(\geq2\) Gross--Zagier open & The needed theorem-grade
higher-rank identity remains an audit-result gap. \\
NC-4 & Additive \(p\leq3\) IMC open & The small-additive-prime input,
including the \(p=2\) part outside the odd-prime carriers, is not proved
here. \\
NC-5 & Imports not eliminated & The imported theorem stacks remain sourced
assumptions. \\
NC-6 & Anchors are not families & The \(\eta\)-formula checks are
anchor-level, not uniform family theorems. \\
NC-7 & AOR non-eliminative & Bridged recognition records are carried, not
zeroed. \\
\bottomrule
\end{tabularx}
\end{table}
  \section{Discussion}\label{sec:discussion}


The point of this paper is not to add another unconditional theorem to
the classical BSD literature.  It is to make explicit what a
conditional closure of Strong BSD rests on, and to organize the
remaining arithmetic content into named records.  The existing
conditional landscape is already rich.  Gross--Zagier--Kolyvagin
methods control the rank \(0\) and \(1\) bridge between derivatives,
Heegner points, Selmer groups, and regulators.  Skinner--Urban and
Iwasawa-main-conjecture methods control ordinary-prime Selmer and
\(p\)-adic \(L\)-function data under their hypotheses.  Cassels--Tate
and Selmer-complex formalisms supply the pairing language needed to
compare arithmetic determinants, Fitting ideals, and local correction
terms.  Each of these classical mechanisms controls part of the
Strong-BSD product.  None, by itself, gives the entire scalar package
in the generality targeted here.

The contribution of the closure is therefore architectural.  Once the
classical theorem stacks are imported, the remaining scalar factors
are named as three recognition sources: the \(p\)-adic descent and
Tamagawa correction, the signed-Selmer persistence of the
\(\Sha\)-factor, and the higher-rank regulator/fixity readout.  The
shell \(\SelBSD\) then records the analytic leading coefficient, the
arithmetic product, and the residual between them in one typed object.
The readout theorem identifies vanishing of that residual with the
standard Strong-BSD scalar identity.  Thus the paper gives a precise
conditional implication: under the standing shell-closure record,
supplying the three recognition sources together with the headline imported
theorem stacks forces Strong BSD.  Its value lies in the bookkeeping: what is imported, what is carried
as recognition content, which factors each source governs, and which
obligations remain open.  This is not a new unconditional proof of
BSD.  It is a conditional closure whose hypotheses have been separated
and named.

This bookkeeping also explains why \citet{TsiokosBSDApparatus2026} is
structurally relevant.  Its typed adequacy and five-column diagnostic
language separates the BSD scalar
package into factors that can be audited independently: height and
regulator data, period data, finite arithmetic data, \(p\)-adic local
data, and determinant-line data.  The \(\kapr\) normalization from
that work is especially important here.  It identifies the cascade
coefficient with the standard factor
\(\#\Sha(E/\Q)/|E(\Q)_{\tors}|^2\), so the higher-rank fixity source
\(\GammaGZ\) can be read as exactly the rank-\(r\) contribution to
the shell residual.  Without this factor-by-factor language, the
conditional implication would be a single opaque scalar assertion.
With it, the closure can attach each remaining obligation to the
factor it governs.

The same structural pattern is shared by the broader Six Birds
Clay-track program.  In each track, the target conjecture is first
encoded as the vanishing or fixed readout of a saturated trace shell.
Deep content that is not derived internally is isolated as named
recognition sources.  The shell readout is then linked to the target
statement by an explicit scalar or structural equivalence, and the
final connect-to-reality step is an Audited Operational Realisability
reading: a refinement-stable, audit-closed membership assertion in the
sense of \cite{TsiokosAOR2026}.  This is the same non-eliminative
``holds in reality'' pattern used in the Riemann-hypothesis,
Navier--Stokes, and P-vs-NP tracks, stated here for Strong BSD.
\Cref{tab:clay-track-parallel} records the parallel in the compact form used
by the series.
\begin{table}[tbp]
\centering
\small
\caption{Structural parallel with the other Six Birds Clay-track papers.}
\label{tab:clay-track-parallel}
\begin{tabularx}{\textwidth}{@{}>{\raggedright\arraybackslash}p{0.16\textwidth}
>{\raggedright\arraybackslash}p{0.26\textwidth}
>{\raggedright\arraybackslash}p{0.28\textwidth}
>{\raggedright\arraybackslash}X@{}}
\toprule
Track & Shell/readout & Recognition content & AOR reading \\
\midrule
Navier--Stokes & Saturated flow shell and regularity readout & analytic
regularity sources & non-eliminative audit closure \\
Riemann hypothesis & Spectral trace shell and zero-location readout &
Selberg-type recognition source & non-eliminative audit closure \\
P-vs-NP & Complexity trace shell and separation readout & formed-computation
recognition sources & non-eliminative audit closure \\
BSD & \(\SelBSD\) and \(\rhoE=0\) readout & \(\GammaPD\), \(\GammaSP\),
\(\GammaGZ\), plus imports & non-eliminative BSD AOR instance \\
\bottomrule
\end{tabularx}
\end{table}

The BSD composite recognition predicate also fits the common
recognition-mode discipline of \cite{TsiokosFoundationsIII2026}.  In
plain terms, the composite carries the same framework trace signature
as the earlier P-vs-NP and Riemann-hypothesis closures: a common
channel-state pattern recording how the composite operator depends on
formed content, source audit, and visible defect elimination.  The
point of noting this here is not to import proofs from the sibling
tracks.  It is to identify a shared structural fingerprint.  BSD is
the third such cross-track instance: its recognition predicate is
computable and falsifiable at the level of the named records, each of
the three recognition sources is load-bearing, and ablation of any one
source leaves a scalar factor of the Strong-BSD package undetermined.

The final AOR reading should therefore be understood carefully.  It
does not eliminate the open arithmetic obligations.  It records that,
under the declared presentation, every residual piece has been either
constructed, imported, or bridged as supplied recognition content.  In
particular, the rank-\(\geq 2\) generalized Gross--Zagier identity and
the additive \(p\leq 3\) Iwasawa main conjecture remain open in the
audited literature.  The closure makes those gaps visible rather than
hiding them.  Its conclusion is an exact conditional landing together
with a non-eliminative audited-realizability reading: if the standing
shell-closure record, the three recognition sources, and the headline
imported theorem stacks are supplied, then the shell residual vanishes,
and the Strong-BSD scalar identity follows.
 \section{Conclusion}\label{sec:conclusion}


This paper constructs the saturated BSD trace shell \(\SelBSD\), whose
residual \(\rhoE\) is the difference between the analytic leading
coefficient and the Strong-BSD arithmetic product.  The readout
equivalence identifies \(\rhoE=0\) with the usual Strong-BSD scalar
identity.  Under the standing shell-closure assumption, and assuming the
three recognition sources---\(p\)-adic descent, signed-Selmer
\(\Sha\)-persistence, and higher-rank fixity---together with the headline
imported theorem stacks---the Cassels--Tate \(\chiCTp\) stack, the
\(\AofE\) Selmer-complex pairing substrate, and the Beilinson
rank-\(\geq 2\) comparison---the conditional landing gives the Strong-BSD
scalar identity.  Outside the framework, the result reads as
\[
  \begin{aligned}
  &\SelBSD\text{-closure}
    \wedge\ \GammaPD \wedge \GammaSP \wedge \GammaGZ\\
  &\quad{}\wedge\ \chiCTp\text{-imports}
    \wedge\ \AofE\text{-imports}
    \wedge\ \text{Beilinson-imports}\\
  &\Longrightarrow \mathrm{Strong\ BSD},
  \end{aligned}
\]
not as an unconditional proof of BSD.

The final AOR reading records the closed carrier as a
refinement-stable, audit-closed instance: informally, the framework's
``Strong BSD holds in reality'' statement under the declared records.
This reading is non-eliminative.  The recognition records are bridged
and carried as supplied arithmetic content; they are not discharged of
their open arithmetic.

The pending obligations remain explicit.  The rank-\(\geq 2\)
generalized Gross--Zagier identity (\(T_{E12}\)) and the additive
\(p\leq 3\) Iwasawa main conjecture (\(T_{\mathrm{BAD}}\)) remain
open audit-result no-gos.  The overconvergent \(\eta\)-formula is
verified only at the stated anchors, not as a uniform family-level
theorem.  Unconditional Strong BSD is therefore deferred.

The contribution is a precise conditional architecture: it names the
arithmetic inputs that close Strong BSD and records the remaining
obligations as explicit, still-open content rather than hiding them in
the scalar conjecture.
 \appendix
\section{Lean formalization}\label{app:formalization}
\subsection*{Trust base and methodology}

The Lean development for this paper is deliberately mathlib-free.  It
does not import an external mathematics library, and therefore no
mathlib axioms occur in the dependency closure of the BSD declarations.
Every declaration entered in the manifest with theorem status was
audited with \texttt{\#print axioms}.  The accepted trust base is
Lean's standard foundational base: \texttt{propext},
\texttt{Quot.sound}, and \texttt{Classical.choice}, with
\texttt{funext} and the \texttt{choice} recursor appearing as standard
consequences in the audit output.  No project-local axiom lies behind
any theorem-marked result.  The source tree also contains no active
\texttt{axiom}, \texttt{opaque}, \texttt{constant}, \texttt{sorry}, or
\texttt{admit} declarations in the BSD namespace.  This point matters
for the interpretation of the recognition sources and imports: they
are represented as typed structure carriers, never as Lean axioms.

The formalization mechanizes the framework-structural closure
architecture, not the imported arithmetic geometry.  There are three
main disclosure modes.  First, typed structure carriers are used for
the imported theorem stacks and the recognition sources.  The
Cassels--Tate, Selmer-complex pairing, Beilinson, overconvergent
\(\eta\), and cascade imports are records whose mathematical content is
carried as hypothesis fields.  The three recognition sources are
likewise records: their arithmetic content is carried explicitly and
then consumed by later statements.  In all of these cases the content
is carried, not derived, and it is not an axiom of the formalization.

Second, several results are projection-packaged conditional schemas.
The \(\chiCTp\) comparison, the overconvergent \(\eta\)-formula, the
rank-\(\leq 1\) cascade theorem, the composite signature, and the
headline conditional landing are checked projections over explicit
proof-carrying records, conditional on the relevant imports and
recognition sources.  The formalization checks that the projections and
implication chains use the supplied records in the stated way; it does
not independently derive the underlying arithmetic theorems.  For this
reason the body footnotes say that these claims are tracked as
conditional schemas, not that the formalization derives the arithmetic
comparison unconditionally.

Third, some statements are substantive finite Lean proofs of the
structural layer.  The readout-level equivalence, the
master-theorem-applicability statement, and the AOR mechanical and
recognition discharge theorems are proofs about the finite record
architecture and its bookkeeping.  The AOR-instance membership is
checked over a narrowed representation: the local syntactic
refinement-stable predicate used by this paper, not the full Audited
Operational Realisability meta-theory.  Thus Lean checks the closure
architecture---the residual readout algebra, the implication chains,
and the discharge bookkeeping---while certifying none of the imported
arithmetic geometry or the recognition-source content itself.

In coverage terms, the paper has 28 statements of record.  Of these,
21 are mechanized: 4 substantive structural theorems, 5
projection-packaged conditional schemas, 8 typed definition carriers
(including the shell and predicate mirrors, the typed import carriers,
and the AOR-instance object), 3 recognition-source carriers, and 1
narrowed AOR-instance membership.  The remaining 7 are not mechanized:
the two audit-result no-gos, two remarks, and three non-claims.  The
tables below give the per-label coverage and the corresponding Lean
declaration map.
\subsection*{Wording-discipline table}

Every mechanized statement in the body carries a footnote whose wording
is canonical, not editorial.  The wording encodes whether the result is
a substantive finite Lean proof, a projection-packaged conditional
schema, a narrowed-representation check, a typed structure carrier, or
a recognition-source carrier.  Projection-packaging and narrowing are
therefore visible at the point of disclosure; unmechanized statements
carry no footnote.

\begin{center}
\begin{tabularx}{\textwidth}{@{}>{\raggedright\arraybackslash}p{0.36\textwidth}X@{}}
\toprule
Statement class & Footnote convention \\
\midrule
Theorem, substantive finite proof &
Verified in Lean as \texttt{[decl]}; see App~\ref{app:formalization}. \\
\addlinespace
Theorem, projection-packaged conditional schema &
Tracked in the formalization harness as a conditional schema
\texttt{[decl]} --- a checked projection over an explicit
proof-carrying record (conditional on the imports and recognition
sources), not an independent Lean derivation. \\
\addlinespace
Theorem, narrowed representation &
Verified in Lean as \texttt{[decl]} over a narrowed representation
(the local syntactic refinement-stable predicate); the narrowing is
recorded below. \\
\addlinespace
Definition, typed structure carrier (import) &
Encoded in Lean as a typed structure carrier \texttt{[decl]} bundling
the imported theorem(s) as hypothesis fields; not a Lean axiom. \\
\addlinespace
Definition, recognition-source carrier &
Supplied as a recognition source, encoded in Lean as a typed structure
carrier \texttt{[decl]} over a narrowed representation; its arithmetic
content is a hypothesis field, not a Lean axiom. \\
\addlinespace
Definition, faithful (typed mirror) &
No footnote; the declaration appears only in the per-label coverage
table below. \\
\addlinespace
No-go / remark / non-claim (not mechanized) &
No footnote and no Lean identifier in the body; recorded as a
paper-level statement only. \\
\bottomrule
\end{tabularx}
\end{center}
\subsection*{Standing narrowings}

\paragraph{Projection-packaged conditional schemas.}
Five results in this paper are mechanized as checked projections over
explicit proof-carrying records rather than as independent Lean
derivations: the \(\chiCTp\) Cassels--Tate comparison, the
overconvergent \(\eta\)-formula, the rank-\(\leq 1\) cascade theorem,
the composite signature, and the headline conditional landing.  This
means that the formalization checks the projection or implication
chain from the supplied records to the stated conclusion.  For example,
in the landing it checks that the recognition sources and headline imports
are threaded through the shell-readout theorems so that the residual
vanishes and the Strong-BSD scalar identity follows.

The corresponding bound is equally important.  These projections do
not independently derive the underlying arithmetic-geometry inputs:
the Cassels--Tate comparison inputs, the overconvergent \(L\)-value
inputs, and the rank-\(\leq 1\) Iwasawa/Heegner inputs remain carried
hypotheses.  This is why the body footnotes for these five results say
that they are tracked as conditional schemas and are not independent
Lean derivations.  In particular, the conditional landing is a theorem
with the recognition sources and headline imports as explicit hypothesis
parameters; it is not an unconditional Lean proof of BSD.

\paragraph{Typed structure carriers.}
The five imported theorem stacks and the three recognition sources are
encoded as typed structures whose arithmetic payload is a hypothesis
field.  Three of the imported stacks, \(\chiCTp\), \(\AofE\), and the
Beilinson comparison, are the headline imports used directly in the landing
theorem; the overconvergent \(\eta\) and \(T_{\mathrm{CASCADE}}\) stacks are
supporting conditional imports recorded by their own sections.  For the
imports, this means that the external theorem stacks are carried as
structured records rather than re-derived.  For the recognition sources, it
means that the named arithmetic content is available as a field of the
carrier and is then consumed by the closure architecture.  In both cases
the content is carried, not derived, and it never appears as a Lean axiom.

The recognition-source carriers also use a narrowed structural
representation.  The carrier records the relevant arithmetic assertion
in typed form---for example the \(p\)-adic descent congruence, the
signed-Selmer persistence comparison, or the higher-rank fixity
readout---without formalizing the full underlying \(p\)-adic, Selmer,
or regulator theory inside Lean.  The dimension and role of each
carrier are faithful to the intended arithmetic content, but the
background arithmetic theory is not redeveloped in the formalization.

Neither of these narrowings enlarges the Lean trust base.  They delimit
what the finite mechanization claims: the imported arithmetic and
recognition content are explicit carried hypotheses, while the checked
Lean layer verifies the structural use of those hypotheses.
\begin{table}[tbp]
\centering
\small
\caption{Standing representation notes for the closure formalization.}
\label{tab:representation-notes}
\begin{tabularx}{\textwidth}{@{}>{\raggedright\arraybackslash}p{0.25\textwidth}
>{\raggedright\arraybackslash}p{0.30\textwidth}
>{\raggedright\arraybackslash}X@{}}
\toprule
Representation issue & Where it appears & Disclosure effect \\
\midrule
Import structures & \(\chiCTp\), \(\AofE\), Beilinson, \(\eta\), and
cascade imports & Arithmetic theorem stacks are hypothesis fields in typed
records, not Lean axioms or internal derivations. \\
\addlinespace
Recognition carriers & \(\GammaPD\), \(\GammaSP\), \(\GammaGZ\) & The
arithmetic payload is supplied closure content and is consumed by the shell
as a hypothesis. \\
\addlinespace
Projection-packaged landing & Supporting conditionals and
\Cref{thm:closure:strong-bsd-conditional} & The formalization checks the
conditional projection/implication chain over supplied records, not an
independent arithmetic proof. \\
\addlinespace
Local \(\RefStable\) predicate & AOR-instance theorem & The Lean theorem
checks the paper-local syntactic predicate, not the full AOR meta-theory. \\
\addlinespace
Cross-paper \(\kapr\) import & \(\Cref{sec:aor-instance}\) & The
normalization is reused from the apparatus paper and cited there, not
re-mechanized in the closure axis. \\
\bottomrule
\end{tabularx}
\end{table}
 \subsection*{Non-eliminative AOR instance}

The AOR-instance membership theorem is mechanized over a narrowed
representation.  The development defines, for this paper, a local
syntactic refinement-stable predicate for the BSD trace shell; it does
not mechanize the full Audited Operational Realisability meta-theory of
\cite{TsiokosAOR2026}.  The checked statement is therefore the following
bounded claim: with respect to this local predicate, the declared audit
records of the BSD trace shell account for the structural defects that
persist under refinement.  It is not a claim that Lean has formalized the
general AOR meta-theory, and it is not a claim of membership in that
full meta-theoretic sense.  This is why the body footnote for the
AOR-instance theorem carries the narrowed-representation qualifier
rather than a bare ``verified in Lean'' formulation.

The non-eliminative content is also explicit in the formal record.  The
three recognition sources are discharged as bridged records: each is
recorded in the audit as a primary source-type structural defect, with
its forced secondary role, target, transport, and limit defects, and is
then reclassified as carried recognition content.  None of the three is
marked as eliminated.  By contrast, the mechanical components---the
trace-shell construction, the \(\kapr\)-normalization, and the imported
theorem stacks---are discharged either by construction or by an approved
external-record citation chain.

Thus the AOR-instance membership is non-eliminative.  It records that
the open arithmetic content has been named, carried, and bridged inside
the audit presentation, not that the content has been closed.  In
particular, the rank-\(\geq 2\) generalized Gross--Zagier identity and
the additive \(p \leq 3\) Iwasawa main conjecture remain open after the
AOR reading; the membership records their place in the declared
presentation rather than proving them.  This subsection fixes the exact
scope of the formal AOR result: a local, refinement-stable,
non-eliminative audit-closed membership, so that neither the headline
conditional landing nor the AOR reading is mistaken for an unconditional
or eliminative theorem.
\subsection*{Per-label coverage}

The table lists every closure statement of record in document order,
together with its environment kind, mechanization fidelity, and body
location.  Rows marked faithful, projection, or narrowed have a Lean
declaration, mapped in the next subsection; rows marked not mechanized
are paper-level statements and carry no Lean declaration.

\begin{small}
\begin{longtable}{@{}>{\ttfamily}p{0.46\textwidth} l l l@{}}
\caption{Per-label coverage for the closure statements of record (28 rows, document order). Labels elide the \texttt{closure:} prefix; ``projection'' marks a projection-packaged conditional schema and ``narrowed'' a narrowed-representation row.}\label{tab:formalization-coverage}\\
\toprule \normalfont Paper label & Kind & Fidelity & Body \\ \midrule \endfirsthead
\multicolumn{4}{@{}l}{\footnotesize\itshape continued from previous page}\\ \toprule \normalfont Paper label & Kind & Fidelity & Body \\ \midrule \endhead
\midrule \multicolumn{4}{r@{}}{\footnotesize\itshape continued on next page}\\ \endfoot
\bottomrule \endlastfoot
def:chi-ct-p-import & \normalfont definition & \normalfont faithful & \normalfont \Cref{sec:imports} \\
def:a-e-import & \normalfont definition & \normalfont faithful & \normalfont \Cref{sec:imports} \\
def:beilinson-import & \normalfont definition & \normalfont faithful & \normalfont \Cref{sec:imports} \\
nonclaim:imports-not-derived & \normalfont nonclaim & \normalfont not mechanized & \normalfont \Cref{sec:imports} \\
def:gamma-padic-descent & \normalfont definition & \normalfont narrowed & \normalfont \Cref{sec:recognition-sources} \\
def:gamma-sha-persistence & \normalfont definition & \normalfont narrowed & \normalfont \Cref{sec:recognition-sources} \\
def:gamma-higher-gz-fixity & \normalfont definition & \normalfont narrowed & \normalfont \Cref{sec:recognition-sources} \\
nonclaim:recognition-sources-not-derived & \normalfont nonclaim & \normalfont not mechanized & \normalfont \Cref{sec:recognition-sources} \\
def:sel-bsd-shell & \normalfont definition & \normalfont faithful & \normalfont \Cref{sec:shell-readout} \\
def:pi-bsd & \normalfont definition & \normalfont faithful & \normalfont \Cref{sec:shell-readout} \\
thm:pi-bsd-iff-strong-bsd & \normalfont theorem & \normalfont faithful & \normalfont \Cref{sec:shell-readout} \\
thm:master-theorem-applicability & \normalfont theorem & \normalfont faithful & \normalfont \Cref{sec:shell-readout} \\
thm:composite-signature & \normalfont theorem & \normalfont projection & \normalfont \Cref{sec:shell-readout} \\
thm:chi-ct-p-comparison & \normalfont theorem & \normalfont projection & \normalfont \Cref{sec:chi-ct-p} \\
def:eta-published-imports & \normalfont definition & \normalfont faithful & \normalfont \Cref{sec:eta-formula} \\
thm:oc-eta-formula & \normalfont theorem & \normalfont projection & \normalfont \Cref{sec:eta-formula} \\
def:t-cascade-import & \normalfont definition & \normalfont faithful & \normalfont \Cref{sec:t-cascade} \\
thm:t-cascade-rank-le-1 & \normalfont theorem & \normalfont projection & \normalfont \Cref{sec:t-cascade} \\
thm:t-e12-no-go & \normalfont theorem & \normalfont not mechanized & \normalfont \Cref{sec:obstructions} \\
thm:t-bad-no-go & \normalfont theorem & \normalfont not mechanized & \normalfont \Cref{sec:obstructions} \\
rmk:mode-b-residuals & \normalfont remark & \normalfont not mechanized & \normalfont \Cref{sec:obstructions} \\
thm:strong-bsd-conditional & \normalfont theorem & \normalfont projection & \normalfont \Cref{sec:landing} \\
nonclaim:landing-conditional & \normalfont nonclaim & \normalfont not mechanized & \normalfont \Cref{sec:landing} \\
def:aor-sel-instance & \normalfont definition & \normalfont faithful & \normalfont \Cref{sec:aor-instance} \\
thm:aor-mechanical-records & \normalfont theorem & \normalfont faithful & \normalfont \Cref{sec:aor-instance} \\
thm:aor-recognition-discharge & \normalfont theorem & \normalfont faithful & \normalfont \Cref{sec:aor-instance} \\
thm:aor-instance & \normalfont theorem & \normalfont narrowed & \normalfont \Cref{sec:aor-instance} \\
rmk:aor-partial-status & \normalfont remark & \normalfont not mechanized & \normalfont \Cref{sec:aor-instance} \\
\end{longtable}
\end{small}
 \subsection*{Lean declaration map}

The 21 mechanized closure statements each correspond to a Lean
declaration, listed here for citation and audit.  Labels elide the
\texttt{closure:} prefix, and declarations elide the
\texttt{SixBirdsBSD.Closure.} namespace; the seven not-mechanized
statements have no declaration and are omitted from this map, with full
coverage recorded in \Cref{tab:formalization-coverage}.  The
\(\kapr\)-normalization reused in the AOR mechanical-records theorem is
imported from \citet{TsiokosBSDApparatus2026};
its own declaration belongs to that paper rather than to the closure
declaration map.

\begin{small}
\begin{longtable}{@{}>{\ttfamily}p{0.44\textwidth} >{\ttfamily}p{0.50\textwidth}@{}}
\caption{Lean declaration map for the 21 mechanized closure statements. Labels elide the \texttt{closure:} prefix; declarations elide the \texttt{SixBirdsBSD.Closure.} namespace.}\label{tab:appd-declmap-closure}\\
\toprule \normalfont Paper label & \normalfont Lean declaration \\ \midrule \endfirsthead
\multicolumn{2}{@{}l}{\footnotesize\itshape continued from previous page}\\ \toprule \normalfont Paper label & \normalfont Lean declaration \\ \midrule \endhead
\midrule \multicolumn{2}{r@{}}{\footnotesize\itshape continued on next page}\\ \endfoot
\bottomrule \endlastfoot
def:\allowbreak chi-\allowbreak ct-\allowbreak p-\allowbreak import & Imports.\allowbreak chi\allowbreak CTp\allowbreak Import \\
def:\allowbreak a-\allowbreak e-\allowbreak import & Imports.\allowbreak a\allowbreak EImport \\
def:\allowbreak beilinson-\allowbreak import & Imports.\allowbreak beilinson\allowbreak Import \\
def:\allowbreak gamma-\allowbreak padic-\allowbreak descent & Recognition\allowbreak Sources.\allowbreak gamma\allowbreak Padic\allowbreak Descent \\
def:\allowbreak gamma-\allowbreak sha-\allowbreak persistence & Recognition\allowbreak Sources.\allowbreak gamma\allowbreak Sha\allowbreak Persistence \\
def:\allowbreak gamma-\allowbreak higher-\allowbreak gz-\allowbreak fixity & Recognition\allowbreak Sources.\allowbreak gamma\allowbreak Higher\allowbreak GZFixity \\
def:\allowbreak sel-\allowbreak bsd-\allowbreak shell & Sel\allowbreak Shell.\allowbreak sel\allowbreak BSDShell \\
def:\allowbreak pi-\allowbreak bsd & Sel\allowbreak Shell.\allowbreak pi\allowbreak BSD \\
thm:\allowbreak pi-\allowbreak bsd-\allowbreak iff-\allowbreak strong-\allowbreak bsd & Sel\allowbreak Shell.\allowbreak pi\allowbreak BSDIff\allowbreak Strong\allowbreak BSD \\
thm:\allowbreak master-\allowbreak theorem-\allowbreak applicability & Sel\allowbreak Shell.\allowbreak master\allowbreak Theorem\allowbreak Applicability \\
thm:\allowbreak composite-\allowbreak signature & Sel\allowbreak Shell.\allowbreak composite\allowbreak Signature \\
thm:\allowbreak chi-\allowbreak ct-\allowbreak p-\allowbreak comparison & Chi\allowbreak CTp.\allowbreak chi\allowbreak CTp\allowbreak Comparison \\
def:\allowbreak eta-\allowbreak published-\allowbreak imports & Eta\allowbreak Formula.\allowbreak eta\allowbreak Published\allowbreak Imports \\
thm:\allowbreak oc-\allowbreak eta-\allowbreak formula & Eta\allowbreak Formula.\allowbreak oc\allowbreak Eta\allowbreak Formula \\
def:\allowbreak t-\allowbreak cascade-\allowbreak import & TCascade.\allowbreak t\allowbreak Cascade\allowbreak Import \\
thm:\allowbreak t-\allowbreak cascade-\allowbreak rank-\allowbreak le-\allowbreak 1 & TCascade.\allowbreak t\allowbreak Cascade\allowbreak Rank\allowbreak Le\allowbreak One \\
thm:\allowbreak strong-\allowbreak bsd-\allowbreak conditional & Landing.\allowbreak strong\allowbreak BSDConditional \\
def:\allowbreak aor-\allowbreak sel-\allowbreak instance & AORInstance.\allowbreak aor\allowbreak Sel\allowbreak Instance \\
thm:\allowbreak aor-\allowbreak mechanical-\allowbreak records & AORInstance.\allowbreak aor\allowbreak Mechanical\allowbreak Records \\
thm:\allowbreak aor-\allowbreak recognition-\allowbreak discharge & AORInstance.\allowbreak aor\allowbreak Recognition\allowbreak Discharge \\
thm:\allowbreak aor-\allowbreak instance & AORInstance.\allowbreak aor\allowbreak Instance \\
\end{longtable}
\end{small}

\clearpage
\begingroup
\raggedright
\small
\bibliographystyle{plainnat}
\begin{thebibliography}{38}
\providecommand{\natexlab}[1]{#1}
\providecommand{\url}[1]{\texttt{#1}}
\expandafter\ifx\csname urlstyle\endcsname\relax
  \providecommand{\doi}[1]{doi: #1}\else
  \providecommand{\doi}{doi: \begingroup \urlstyle{rm}\Url}\fi

\bibitem[Beilinson(1985)]{Beilinson1985}
A.~A. Beilinson.
\newblock Higher regulators and values of {L}-functions.
\newblock \emph{J. Soviet Math.}, 30:\penalty0 2036--2070, 1985.
\newblock \doi{10.1007/BF02105861}.

\bibitem[Bella\"iche(2012)]{Bellaiche2012}
Jo\"el Bella\"iche.
\newblock Critical p-adic {L}-functions.
\newblock \emph{Invent. Math.}, 189\penalty0 (1):\penalty0 1--60, 2012.
\newblock \doi{10.1007/s00222-011-0358-z}.

\bibitem[Benois and B\"uy\"ukboduk(2024)]{BenoisBuyukboduk2024}
Denis Benois and K\^az\i~m B\"uy\"ukboduk.
\newblock Arithmetic of critical p-adic {L}-functions, 2024.
\newblock arXiv:2403.16076.

\bibitem[Birch and Swinnerton-Dyer(1965)]{BirchSwinnertonDyer1965}
B.~J. Birch and H.~P.~F. Swinnerton-Dyer.
\newblock Notes on elliptic curves. {II}.
\newblock \emph{J. Reine Angew. Math.}, 218:\penalty0 79--108, 1965.
\newblock \doi{10.1515/crll.1965.218.79}.

\bibitem[Bloch and Kato(1990)]{BlochKato1990}
Spencer Bloch and Kazuya Kato.
\newblock {L}-functions and {T}amagawa numbers of motives.
\newblock In \emph{The {Grothendieck} {Festschrift}, Vol.~I}, volume~86 of \emph{Progr. Math.}, pages 333--400. Birkh\"auser, 1990.
\newblock \doi{10.1007/978-0-8176-4574-8_9}.

\bibitem[Burns and Flach(2001)]{BurnsFlach2001}
David Burns and Matthias Flach.
\newblock Tamagawa numbers for motives with (non-commutative) coefficients.
\newblock \emph{Doc. Math.}, 6:\penalty0 501--570, 2001.
\newblock \doi{10.4171/DM/113}.

\bibitem[Burns et~al.(2024)Burns, Macias~Castillo, and Seo]{BurnsMaciasCastilloSeo2024}
David Burns, Daniel Macias~Castillo, and Soogil Seo.
\newblock On {Weil--Stark} elements, {I}: construction and general properties.
\newblock \emph{J. Lond. Math. Soc. (2)}, 110\penalty0 (2):\penalty0 e70001, 2024.
\newblock \doi{10.1112/jlms.70001}.

\bibitem[Burungale et~al.(2022)Burungale, Castella, Skinner, and Tian]{BurungaleCastellaSkinnerTian2022}
Ashay Burungale, Francesc Castella, Christopher Skinner, and Ye~Tian.
\newblock {$p^\infty$}-{Selmer} groups and rational points on {CM} elliptic curves.
\newblock \emph{Ann. Math. Qu\'e.}, 46\penalty0 (2):\penalty0 325--346, 2022.
\newblock \doi{10.1007/s40316-022-00203-y}.

\bibitem[Cassels(1962)]{Cassels1962}
J.~W.~S. Cassels.
\newblock Arithmetic on curves of genus 1. {IV}. {Proof} of the {Hauptvermutung}.
\newblock \emph{J. Reine Angew. Math.}, 211:\penalty0 95--112, 1962.

\bibitem[Castella et~al.(2025)Castella, Hsu, Kundu, Lee, and Liu]{CHKLL2025}
Francesc Castella, Chi-Yun Hsu, Debanjana Kundu, Yu-Sheng Lee, and Zheng Liu.
\newblock Derived p-adic heights and the leading coefficient of the {Bertolini--Darmon--Prasanna} p-adic {L}-function.
\newblock \emph{Trans. Amer. Math. Soc. Ser. B}, 12:\penalty0 748--788, 2025.
\newblock \doi{10.1090/btran/227}.

\bibitem[Coleman and Mazur(1998)]{ColemanMazur1998}
Robert Coleman and Barry Mazur.
\newblock The eigencurve.
\newblock In \emph{Galois Representations in Arithmetic Algebraic Geometry ({Durham}, 1996)}, volume 254 of \emph{London Math. Soc. Lecture Note Ser.}, pages 1--113. Cambridge Univ. Press, 1998.

\bibitem[Fisher et~al.(2010)Fisher, Schaefer, and Stoll]{FisherSchaeferStoll2010}
Tom Fisher, Edward~F. Schaefer, and Michael Stoll.
\newblock The yoga of the {Cassels--Tate} pairing.
\newblock \emph{LMS J. Comput. Math.}, 13:\penalty0 451--460, 2010.
\newblock \doi{10.1112/S1461157010000185}.

\bibitem[Flach(1990)]{Flach1990}
Matthias Flach.
\newblock A generalisation of the {Cassels--Tate} pairing.
\newblock \emph{J. Reine Angew. Math.}, 412:\penalty0 113--127, 1990.
\newblock \doi{10.1515/crll.1990.412.113}.

\bibitem[Fontaine and Perrin-Riou(1994)]{FontainePerrinRiou1994}
Jean-Marc Fontaine and Bernadette Perrin-Riou.
\newblock Autour des conjectures de {Bloch} et {Kato}: cohomologie galoisienne et valeurs de fonctions {L}.
\newblock In \emph{Motives ({Seattle}, 1991)}, volume 55.1 of \emph{Proc. Sympos. Pure Math.}, pages 599--706. Amer. Math. Soc., 1994.

\bibitem[Gross and Zagier(1986)]{GrossZagier1986}
Benedict~H. Gross and Don~B. Zagier.
\newblock Heegner points and derivatives of {L}-series.
\newblock \emph{Invent. Math.}, 84\penalty0 (2):\penalty0 225--320, 1986.
\newblock \doi{10.1007/BF01388809}.

\bibitem[Knudsen and Mumford(1976)]{KnudsenMumford1976}
Finn~Faye Knudsen and David Mumford.
\newblock The projectivity of the moduli space of stable curves. {I}. {Preliminaries} on ``det'' and ``div''.
\newblock \emph{Math. Scand.}, 39\penalty0 (1):\penalty0 19--55, 1976.
\newblock \doi{10.7146/math.scand.a-11642}.

\bibitem[Kobayashi(2003)]{Kobayashi2003}
Shin-ichi Kobayashi.
\newblock Iwasawa theory for elliptic curves at supersingular primes.
\newblock \emph{Invent. Math.}, 152\penalty0 (1):\penalty0 1--36, 2003.
\newblock \doi{10.1007/s00222-002-0265-4}.

\bibitem[Kolyvagin(1989)]{Kolyvagin1989}
V.~A. Kolyvagin.
\newblock Finiteness of {$E(\mathbb{Q})$} and {$\Sha(E,\mathbb{Q})$} for a subclass of {Weil} curves.
\newblock \emph{Math. USSR-Izv.}, 32\penalty0 (3):\penalty0 523--541, 1989.
\newblock \doi{10.1070/IM1989v032n03ABEH000779}.

\bibitem[Lang and Wake(2025)]{LangWake2025}
Jaclyn Lang and Preston Wake.
\newblock The {Eisenstein} ideal at prime-square level has constant rank.
\newblock \emph{Proc. Natl. Acad. Sci. USA}, 122\penalty0 (28):\penalty0 e2500729122, 2025.
\newblock \doi{10.1073/pnas.2500729122}.

\bibitem[Ling(1997)]{Ling1997}
San Ling.
\newblock On the {$\mathbb{Q}$}-rational cuspidal subgroup and the component group of {$J_0(p^r)$}.
\newblock \emph{Israel J. Math.}, 99:\penalty0 29--54, 1997.
\newblock \doi{10.1007/BF02760675}.

\bibitem[Lorenzini(1993)]{Lorenzini1993}
Dino~J. Lorenzini.
\newblock On the group of components of a {N\'eron} model.
\newblock \emph{J. Reine Angew. Math.}, 445:\penalty0 109--160, 1993.
\newblock \doi{10.1515/crll.1993.445.109}.

\bibitem[Macias~Castillo and Sano(2026)]{MaciasCastilloSano2026}
Daniel Macias~Castillo and Takamichi Sano.
\newblock On {Selmer} complexes, {Stark} systems and derived p-adic heights, 2026.
\newblock Preprint, arXiv:2603.23978.

\bibitem[Mazur and Rubin(2004)]{MazurRubin2004}
Barry Mazur and Karl Rubin.
\newblock Kolyvagin systems.
\newblock \emph{Mem. Amer. Math. Soc.}, 168\penalty0 (799), 2004.
\newblock \doi{10.1090/memo/0799}.

\bibitem[Nekov\'a\v{r}(2006)]{Nekovar2006}
Jan Nekov\'a\v{r}.
\newblock \emph{Selmer Complexes}.
\newblock Number 310 in Ast\'erisque. Soc. Math. France, 2006.

\bibitem[Pollack(2003)]{Pollack2003}
Robert Pollack.
\newblock On the p-adic {L}-function of a modular form at a supersingular prime.
\newblock \emph{Duke Math. J.}, 118\penalty0 (3):\penalty0 523--558, 2003.
\newblock \doi{10.1215/S0012-7094-03-11835-9}.

\bibitem[Pollack and Stevens(2011)]{PollackStevens2011}
Robert Pollack and Glenn Stevens.
\newblock Overconvergent modular symbols and p-adic {L}-functions.
\newblock \emph{Ann. Sci. \'Ec. Norm. Sup\'er. (4)}, 44\penalty0 (1):\penalty0 1--42, 2011.
\newblock \doi{10.24033/asens.2139}.

\bibitem[Poonen and Stoll(1999)]{PoonenStoll1999}
Bjorn Poonen and Michael Stoll.
\newblock The {Cassels--Tate} pairing on polarized abelian varieties.
\newblock \emph{Ann. of Math. (2)}, 150\penalty0 (3):\penalty0 1109--1149, 1999.
\newblock \doi{10.2307/121064}.

\bibitem[Rubin(1991)]{Rubin1991}
Karl Rubin.
\newblock The ``main conjectures'' of {Iwasawa} theory for imaginary quadratic fields.
\newblock \emph{Invent. Math.}, 103\penalty0 (1):\penalty0 25--68, 1991.
\newblock \doi{10.1007/BF01239508}.

\bibitem[Sakamoto(2018)]{Sakamoto2018}
Ryotaro Sakamoto.
\newblock Stark systems over {Gorenstein} local rings.
\newblock \emph{Algebra Number Theory}, 12\penalty0 (10):\penalty0 2295--2326, 2018.
\newblock \doi{10.2140/ant.2018.12.2295}.

\bibitem[Serre(1972)]{Serre1972}
Jean-Pierre Serre.
\newblock Propri\'et\'es galoisiennes des points d'ordre fini des courbes elliptiques.
\newblock \emph{Invent. Math.}, 15\penalty0 (4):\penalty0 259--331, 1972.
\newblock \doi{10.1007/BF01405086}.

\bibitem[Silverman(1994)]{Silverman1994}
Joseph~H. Silverman.
\newblock \emph{Advanced Topics in the Arithmetic of Elliptic Curves}, volume 151 of \emph{Graduate Texts in Mathematics}.
\newblock Springer, 1994.
\newblock \doi{10.1007/978-1-4612-0851-8}.

\bibitem[Skinner and Urban(2014)]{SkinnerUrban2014}
Christopher Skinner and Eric Urban.
\newblock The {Iwasawa} main conjectures for {$\mathrm{GL}_2$}.
\newblock \emph{Invent. Math.}, 195\penalty0 (1):\penalty0 1--277, 2014.
\newblock \doi{10.1007/s00222-013-0448-1}.

\bibitem[Tate(1966)]{Tate1966bsd}
John~T. Tate.
\newblock On the conjectures of {Birch} and {Swinnerton-Dyer} and a geometric analog.
\newblock In \emph{S\'eminaire Bourbaki, Vol.~9}, pages Exp.~No.~306, 415--440. Soc. Math. France, 1966.

\bibitem[Tsiokos(2026{\natexlab{a}})]{TsiokosAOR2026}
Ioannis Tsiokos.
\newblock {Audited Operational Realisability and the Closure Completion of PreSB}, 2026{\natexlab{a}}.

\bibitem[Tsiokos(2026{\natexlab{b}})]{TsiokosBSDApparatus2026}
Ioannis Tsiokos.
\newblock {Decomposing the Bloch--Kato Fundamental Line: Adequacy and No-Go Results for the BSD Invariants}, 2026{\natexlab{b}}.
\newblock Preprint. \doi{10.5281/zenodo.20713981}.

\bibitem[Tsiokos(2026{\natexlab{c}})]{TsiokosFoundationsIII2026}
Ioannis Tsiokos.
\newblock {Six Birds Foundations III: A Finite Audited Interaction Calculus for SBT}, 2026{\natexlab{c}}.

\bibitem[Yan and Zhu(2026)]{YanZhu2026}
Xiaojun Yan and Xiuwu Zhu.
\newblock Main conjectures for non-{CM} elliptic curves at good ordinary primes.
\newblock \emph{J. Algebra}, 693:\penalty0 372--402, 2026.
\newblock \doi{10.1016/j.jalgebra.2026.01.016}.

\bibitem[Zhang(2014)]{WeiZhang2014}
Wei Zhang.
\newblock Selmer groups and the indivisibility of {Heegner} points.
\newblock \emph{Camb. J. Math.}, 2\penalty0 (2):\penalty0 191--253, 2014.
\newblock \doi{10.4310/CJM.2014.v2.n2.a2}.

\end{thebibliography}

\endgroup
\end{document}
